\unit{Bài trình bày 190 — Kỹ thuật hạ xuống I}
\sid{190}
\markboth{Thuyết trình 190}{Thuyết trình 190}
\src{299}
\seminar{Năm thứ 12, 1959/60, số 190\\Tháng 12 năm 1959}
\paperhead{KỸ THUẬT HẠ XUỐNG VÀ CÁC ĐỊNH LÝ TỒN TẠI TRONG HÌNH HỌC ĐẠI SỐ\\
I. ĐẠI CƯƠNG. HẠ XUỐNG THEO CÁC CẤU XẠ PHẲNG TRUNG THÀNH}%
  {do Alexander GROTHENDIECK}

Xét về mặt kỹ thuật, bài trình bày này và các bài tiếp theo có thể được xem là
những biến thể của ``định lý 90'' nổi tiếng của Hilbert. Có vẻ như A.~Weil là
người đưa phương pháp hạ xuống vào hình học đại số, dưới tên gọi ``hạ xuống
trường cơ sở''. Tuy nhiên, Weil chỉ xét trường hợp các mở rộng trường hữu hạn
tách được. Sau đó P.~Cartier xử lý trường hợp các mở rộng thuần bất khả phân có
chiều cao 1. Vì chưa có ngôn ngữ lược đồ, và cụ thể hơn là vì chưa cho phép xét
các phần tử lũy linh trong các vành, Cartier đã không thể phát biểu sự đồng nhất
cốt yếu giữa hai trường hợp này.

Hiện nay, có vẻ như kỹ thuật hạ xuống tổng quát sẽ được trình bày (kết hợp, khi
thích hợp, với các định lý cơ bản của ``hình học hình thức'', xem [3]) là nền
tảng của phần lớn các định lý tồn tại trong hình học đại số. Cũng cần lưu ý rằng
kỹ thuật ấy chắc chắn có thể được chuyển sang ``hình học giải tích'', và có thể
hy vọng rằng, trong tương lai gần, các chuyên gia sẽ chứng minh được những định
lý tương tự theo nghĩa ``giải tích'' của các định lý tồn tại trong hình học
hình thức được nêu ở Bài trình bày II. Dù sao, các công trình gần đây của Kodaira--Spencer,
với những phương pháp dường như không thích hợp để định nghĩa và nghiên cứu các
``đa tạp môđuli'' ở lân cận các điểm kỳ dị của chúng, cho thấy khá rõ nhu cầu về
những phương pháp gần với lý thuyết lược đồ hơn (dĩ nhiên các kỹ thuật siêu việt
sẽ phải bổ sung cho chúng).

Trong Bài trình bày I này, ta sẽ xét trường hợp hạ xuống sơ đẳng nhất (được nêu
trong tiêu đề). Tuy vậy, các ứng dụng của các Định lý 1, 2 và 3 dưới đây đã rất
phong phú. Ta chỉ nêu một vài ứng dụng làm ví dụ, mà không cố phát biểu chúng ở
mức tổng quát tối đa có thể có tại đây.

Ta sẽ tự do sử dụng ngôn ngữ lược đồ; về ngôn ngữ này, xin xem bài trình bày đã
dẫn và [2]. Tuy nhiên, cần lưu ý rõ rằng các tiền lược đồ được xét trong bài này
không nhất thiết là Noether, và các

\src{300}
cấu xạ không nhất thiết thuộc kiểu hữu hạn. Chẳng hạn, nếu $A$ là một vành địa
phương Noether có vành đầy đủ hóa $\widehat A$, ta sẽ phải xét vành không
Noether $\widehat A\otimes_A\widehat A$, cùng các cấu xạ giữa những lược đồ
affine tương ứng với các phép nhúng giữa các vành đang xét.

\fsection{A}{Các kiến thức chuẩn bị về phạm trù}

\subsection*{1. Phạm trù phân thớ, dữ liệu hạ xuống, cấu xạ
$\mathcal F$-hạ xuống}
\addcontentsline{toc}{subsection}{1. Phạm trù phân thớ và dữ liệu hạ xuống}

\result{a. ĐỊNH NGHĨA}{1.1} Một \emph{phạm trù phân thớ} $\mathcal F$ trên cơ
sở $\mathcal C$ gồm một phạm trù $\mathcal C$; với mỗi
$X\in\mathcal C$, một phạm trù $\mathcal F_X$; với mỗi cấu xạ trong
$\mathcal C$, $f:X\longrightarrow Y$, một hàm tử
$f^*:\mathcal F_Y\longrightarrow\mathcal F_X$, cũng được ký hiệu bởi
\[
  f^*(\xi)=\xi\times_Y X
\]
với $\xi\in\mathcal F_Y$ (trong đó $X$ được xem là một ``đối tượng của
$\mathcal C$ trên $Y$'', nghĩa là được trang bị cấu xạ $f$); và cuối cùng, với
hai cấu xạ hợp thành được
$X\xrightarrow{f}Y\xrightarrow{g}Z$, một đẳng cấu hàm tử
\[
  c_{f,g}:(gf)^*\longrightarrow f^*g^*,
\]
sao cho các dữ liệu này thỏa mãn những điều kiện sau:

\begin{enumerate}
\item[(i)] $\operatorname{id}_X^*=\operatorname{id}$;
\item[(ii)] $c_{f,g}$ là đẳng cấu đồng nhất nếu $f$ hoặc $g$ là một đẳng cấu
đồng nhất;
\item[(iii)] với ba cấu xạ hợp thành được
$X\xrightarrow{f}Y\xrightarrow{g}Z\xrightarrow{h}T$, biểu đồ các đẳng cấu sau,
được tạo bởi những đẳng cấu dạng $c_{u,v}$, là giao hoán:
\[
\begin{array}{ccc}
  (h(gf))^*=((hg)f)^* & \longrightarrow & f^*(hg)^* \\
  \big\downarrow && \big\downarrow \\
  (gf)^*h^* & \longrightarrow & (f^*g^*)h^*=f^*(g^*h^*).
\end{array}
\]
\end{enumerate}

\result{VÍ DỤ}{1} Cho $\mathcal C$ là một phạm trù trong đó tồn tại các tích
thớ. Khi đó ta định nghĩa một phạm trù phân thớ $\mathcal F$ trên cơ sở
$\mathcal C$ bằng cách lấy $\mathcal F_X=$ phạm trù các đối tượng của
$\mathcal C$ trên $X$; hàm tử
$f^*:\mathcal F_Y\longrightarrow\mathcal F_X$ ứng với một cấu xạ
$f:X\longrightarrow Y$ được định nghĩa bằng tích thớ
$Z\longmapsto Z\times_YX$.

\result{VÍ DỤ}{2} Cho $\mathcal C$ là phạm trù các tiền lược đồ và, với
$X\in\mathcal C$, cho $\mathcal F_X$ là phạm trù các bó môđun tựa nhất quán
trên $X$. Nếu $f:X\longrightarrow Y$ là một cấu xạ tiền lược đồ, thì
$f^*:\mathcal F_Y\longrightarrow\mathcal F_X$ là hàm tử ảnh ngược

\src{301}
trên các bó môđun. Như vậy ta thu được một phạm trù phân thớ trên cơ sở
$\mathcal C$.

\result{b. ĐỊNH NGHĨA}{1.2} Một biểu đồ
\[
  E\xrightarrow{u}E'
  \overset{v_1}{\underset{v_2}{\rightrightarrows}}E''
\]
gồm các ánh xạ giữa những tập hợp được gọi là \emph{khớp} nếu $u$ là một song
ánh từ $E$ lên phần của $E'$ gồm các $x'\in E'$ thỏa mãn
$v_1(x')=v_2(x')$.

\result{ĐỊNH NGHĨA}{1.3} Cho $\mathcal F$ là một phạm trù phân thớ trên cơ sở
$\mathcal C$; xét một biểu đồ
\[
  \tag{+}
  S\xleftarrow{\alpha}S'
  \overset{\beta_1}{\underset{\beta_2}{\rightrightarrows}}S''
\]
gồm các cấu xạ trong $\mathcal C$ sao cho
$\alpha\beta_1=\alpha\beta_2$; biểu đồ này được gọi là
$\mathcal F$-\emph{khớp} nếu, với mọi cặp đối tượng $\xi,\eta$ của
$\mathcal F_S$, biểu đồ các ánh xạ giữa tập hợp sau là khớp:
\[
  \tag{+'}
  \operatorname{Hom}(\xi,\eta)
  \xrightarrow{\alpha^*}
  \operatorname{Hom}(\alpha^*\xi,\alpha^*\eta)
  \overset{\beta_1^*}{\underset{\beta_2^*}{\rightrightarrows}}
  \operatorname{Hom}(\gamma^*\xi,\gamma^*\eta),
\]
trong đó $\gamma=\alpha\beta_1=\alpha\beta_2$. Để đơn giản, trong biểu đồ cuối
ta đã đồng nhất $\beta_i^*\alpha^*$ với
$(\alpha\beta_i)^*=\gamma^*$ nhờ $c_{\beta_i,\alpha}$.

\result{ĐỊNH NGHĨA}{1.4} Cho $\mathcal F$ là một phạm trù phân thớ trên cơ sở
$\mathcal C$; xét hai cấu xạ $\beta_1,\beta_2$ từ $S''$ đến $S'$ trong
$\mathcal C$. Cho $\xi'\in\mathcal F_{S'}$. Một \emph{dữ liệu dán} trên
$\xi'$ (đối với cặp $(\beta_1,\beta_2)$) là một đẳng cấu từ
$\beta_1^*(\xi')$ lên $\beta_2^*(\xi')$. Nếu
$\xi',\eta'\in\mathcal F_{S'}$ đều được trang bị một dữ liệu dán, thì một cấu
xạ $u:\xi'\longrightarrow\eta'$ trong $\mathcal F_{S'}$ được gọi là
\emph{tương thích với các dữ liệu dán} nếu biểu đồ sau là giao hoán:
\[
\begin{array}{ccc}
  \beta_1^*(\xi') & \longrightarrow & \beta_2^*(\xi') \\
  \big\downarrow && \big\downarrow \\
  \beta_1^*(\eta') & \longrightarrow & \beta_2^*(\eta').
\end{array}
\]
Theo cách này, các đối tượng của $\mathcal F_{S'}$ được trang bị một dữ liệu
dán (đối với $\beta_1,\beta_2$) tạo thành một phạm trù. Nếu
$\alpha:S'\longrightarrow S$ là một cấu xạ sao cho
$\alpha\beta_1=\alpha\beta_2$, thì với mọi $\xi\in\mathcal F_S$, đối tượng
$\xi'=\alpha^*(\xi)$

\src{302}
của $\mathcal F_{S'}$ được trang bị một dữ liệu dán chính tắc, bởi vì
\[
  \beta_i^*\alpha^*(\xi)\simeq(\alpha\beta_i)^*(\xi)=\gamma^*(\xi),
\]
trong đó một lần nữa $\gamma=\alpha\beta_1=\alpha\beta_2$; hơn nữa, nếu
$u:\xi\longrightarrow\eta$ là một cấu xạ trong $\mathcal F_S$, thì
\[
  \alpha^*(u):\alpha^*(\xi)\longrightarrow\alpha^*(\eta)
\]
là một cấu xạ trong $\mathcal F_{S'}$ tương thích với các dữ liệu dán chính
tắc. Do đó ta thu được một hàm tử chính tắc từ phạm trù $\mathcal F_S$ đến
phạm trù các đối tượng của $\mathcal F_{S'}$ được trang bị một dữ liệu dán đối
với cặp $(\beta_1,\beta_2)$. Nói cách khác, Định nghĩa 1.3 cũng có thể được
phát biểu rằng biểu đồ (+) trong định nghĩa ấy là $\mathcal F$-khớp nếu hàm tử
trên là ``đầy đủ và trung thành'', nghĩa là nó xác định một tương đương giữa
phạm trù $\mathcal F_S$ và một phạm trù con của phạm trù các đối tượng
của $\mathcal F_{S'}$ được trang bị dữ liệu dán đối với
$(\beta_1,\beta_2)$.

\result{ĐỊNH NGHĨA}{1.5} Một dữ liệu dán trên
$\xi'\in\mathcal F_{S'}$ được gọi là \emph{hữu hiệu} (đối với $\alpha$) nếu
$\xi'$, cùng dữ liệu này, đẳng cấu với một $\alpha^*(\xi)$ nào đó, trong đó
$\xi\in\mathcal F_S$.

Khi biểu đồ (+) là $\mathcal F$-khớp, đối tượng $\xi$ trong định nghĩa trên vì
thế được xác định tới một đẳng cấu duy nhất, và phạm trù
$\mathcal F_S$ tương đương với phạm trù các đối tượng của
$\mathcal F_{S'}$ được trang bị một dữ liệu dán hữu hiệu.

\medskip\noindent
c. Trường hợp quan trọng nhất là khi
\[
  S''=S'\times_S S',
\]
trong đó các $\beta_i$ là hai phép chiếu $p_1$ và $p_2$ từ
$S'\times_S S'$ lên hai thừa số của nó (bây giờ ta giả sử rằng các tích thớ tồn
tại trong $\mathcal C$). Khi đó có ngay hai điều kiện cần để một dữ liệu dán
$\varphi:p_1^*(\xi')\longrightarrow p_2^*(\xi')$ trên một
$\xi'\in\mathcal F_{S'}$ là hữu hiệu. Ta có
\[
  \tag{i}
  \Delta^*(\varphi)=\operatorname{id}_{\xi'},
\]
trong đó $\Delta:S'\longrightarrow S'\times_S S'$ là cấu xạ đường chéo và ta
đã đồng nhất $\Delta^*p_i^*(\xi')$ với
$(p_i\Delta)^*(\xi')=\xi'$. Ta cũng có
\[
  \tag{ii}
  p_{31}^*(\varphi)=p_{32}^*(\varphi)\,p_{21}^*(\varphi),
\]

\src{303}
trong đó $p_{ij}$ là phép chiếu chính tắc từ
$S'\times_S S'\times_S S'$ lên tích riêng phần của thừa số thứ $i$ và thừa số
thứ $j$.

\result{ĐỊNH NGHĨA}{1.6} Một \emph{dữ liệu hạ xuống} trên
$\xi'\in\mathcal F_{S'}$, đối với cấu xạ
$\alpha:S'\longrightarrow S$, là một dữ liệu dán trên $\xi'$ đối với cặp
$(p_1,p_2)$ các phép chiếu chính tắc
$S'\times_S S'\longrightarrow S'$ thỏa mãn các điều kiện (i) và (ii) ở trên.

\result{ĐỊNH NGHĨA}{1.7} Một cấu xạ $\alpha:S'\longrightarrow S$ được gọi là
một \emph{cấu xạ $\mathcal F$-hạ xuống} nếu biểu đồ các cấu xạ
\[
  S\xleftarrow{\alpha}S'
  \overset{p_1}{\underset{p_2}{\rightrightarrows}}S'\times_S S'
\]
là $\mathcal F$-khớp (Định nghĩa 1.3); $\alpha$ được gọi là một
\emph{cấu xạ $\mathcal F$-hạ xuống nghiêm ngặt} nếu thêm vào đó, mọi dữ liệu hạ
xuống (Định nghĩa 1.6) trên một đối tượng của $\mathcal F_{S'}$ đều hữu hiệu.

Điều kiện cuối cũng có thể được diễn đạt trực quan hơn như sau: việc cho một
đối tượng của $\mathcal F_S$ là ``tương đương'' với việc cho một đối tượng của
$\mathcal F_{S'}$ được trang bị dữ liệu hạ xuống.

Lưu ý rằng nếu một cấu xạ $\alpha:S'\longrightarrow S$ có một ``tiết diện''
$s:S\longrightarrow S'$ (nghĩa là một cấu xạ $s$ sao cho
$\alpha s=\operatorname{id}_S$), thì nó là một cấu xạ
$\mathcal F$-hạ xuống nghiêm ngặt. Thật vậy, nếu
$\xi'\in\mathcal F_{S'}$ được trang bị một dữ liệu hạ xuống đối với $\alpha$,
thì nó ``đến từ'' $\xi=s^*(\xi')$.

\medskip\noindent
d. Có thể trình bày các khái niệm trên một cách trực quan hơn bằng cách đưa vào,
với một đối tượng tham số $T$ của $\mathcal C$ trên $S$, tập hợp
\[
  \operatorname{Hom}_S(T,S')=S'(T),
\]
có các phần tử được ký hiệu bởi $t,t'$, v.v. Với một đối tượng
$\xi'\in\mathcal F_{S'}$ đã cho, mỗi $t\in S'(T)$ khi đó tương ứng với một đối
tượng $t^*(\xi')$ của $\mathcal F_T$, cũng được ký hiệu là $\xi'_t$. Một dữ
liệu dán trên $\xi'$ (đối với $p_1,p_2$) khi đó được xác định bằng cách cho,
với mọi $T$ trên $S$ và mọi cặp điểm $t,t'\in S'(T)$, một đẳng cấu
\[
  \varphi_{t',t}:\xi'_t\longrightarrow\xi'_{t'}
\]
(thỏa mãn các điều kiện hiển nhiên về tính hàm tử khi $T$ biến thiên). Các điều
kiện (i) và (ii) của 1.c khi đó được viết thành
\[
  \tag{i bis}
  \varphi_{t,t}=\operatorname{id}_{\xi'_t}
  \qquad\text{với mọi $T$ và mọi $t\in S'(T)$},
\]

\src{304}
\[
  \tag{ii bis}
  \varphi_{t,t''}=\varphi_{t,t'}\varphi_{t',t''}
  \qquad\text{với mọi $T$ và mọi $t,t',t''\in S'(T)$}.
\]
Hơn nữa, lấy $t=t'=t''$ trong (ii bis), ta thấy
$\varphi_{t,t}=\varphi_{t,t}^2$; vì theo giả thiết $\varphi_{t,t}$ là một đẳng
cấu, suy ra hệ thức (i bis), nên thật ra (i bis) là hệ quả của (ii bis) (do đó
(i) là hệ quả của (ii)). Nhưng nếu không còn giả sử tiên nghiệm rằng các
$\varphi_{t,t'}$ là đẳng cấu (nghĩa là
$\varphi:p_1^*(\xi')\longrightarrow p_2^*(\xi')$ là một đẳng cấu), thì (ii bis)
không nhất thiết kéo theo (i bis); tuy vậy, kết hợp (ii bis) với (i bis) suy ra
rằng các $\varphi_{t,t'}$ là đẳng cấu (bởi ta có
$\varphi_{t,t'}\varphi_{t',t}=\varphi_{t,t}=\operatorname{id}_{\xi'_t}$).

\subsection*{2. Các biểu đồ khớp và các toàn cấu nghiêm ngặt, các cấu xạ
hạ xuống. Ví dụ}
\addcontentsline{toc}{subsection}{2. Các biểu đồ khớp và các toàn cấu nghiêm ngặt}

\result{a. ĐỊNH NGHĨA}{2.1} Cho $\mathcal C$ là một phạm trù. Một biểu đồ
\[
  T\xrightarrow{\alpha}T'
  \overset{\beta_1}{\underset{\beta_2}{\rightrightarrows}}T''
\]
gồm các cấu xạ trong $\mathcal C$ được gọi là \emph{khớp} nếu, với mọi
$Z\in\mathcal C$, biểu đồ tương ứng gồm các ánh xạ giữa tập hợp
\[
  \operatorname{Hom}(Z,T)\longrightarrow
  \operatorname{Hom}(Z,T')
  \rightrightarrows\operatorname{Hom}(Z,T'')
\]
là khớp (Định nghĩa 1.2). Khi đó ta nói $(T,\alpha)$ (hoặc, theo lối nói lạm
dụng, $T$) là một \emph{hạt nhân} của cặp cấu xạ
$(\beta_1,\beta_2)$.

Hạt nhân này hiển nhiên được xác định tới một đẳng cấu duy nhất. Nếu
$\mathcal C$ là phạm trù các tập hợp, thì định nghĩa trên tương thích với Định
nghĩa 1.2. Theo cách đối ngẫu, ta định nghĩa tính khớp của một biểu đồ
\[
  S\xleftarrow{\alpha}S'
  \overset{\beta_1}{\underset{\beta_2}{\rightrightarrows}}S''
\]
gồm các cấu xạ trong $\mathcal C$; khi đó ta nói $(S,\alpha)$ là một
\emph{đối hạt nhân} của cặp cấu xạ $(\beta_1,\beta_2)$.

\result{ĐỊNH NGHĨA}{2.2} Một cấu xạ $\alpha:S'\longrightarrow S$ được gọi là
một \emph{toàn cấu nghiêm ngặt} nếu nó là một toàn cấu và nếu, với mọi cấu xạ
$u:S'\longrightarrow Z$, điều kiện cần sau đây cũng đủ để $u$ phân tích thành
$S'\longrightarrow S\longrightarrow Z$: với mọi $S''\in\mathcal C$ và mọi
cặp cấu xạ $\beta_1,\beta_2:S''\longrightarrow S'$ sao cho
$\alpha\beta_1=\alpha\beta_2$, ta cũng có $u\beta_1=u\beta_2$.

Nếu tích thớ $S'\times_S S'$ tồn tại, điều này tương đương với việc biểu đồ
\[
  S\xleftarrow{\alpha}S'
  \overset{p_1}{\underset{p_2}{\rightrightarrows}}S'\times_S S'
\]

\src{305}
là khớp, nghĩa là $S$ là một đối hạt nhân của cặp $(p_1,p_2)$. Trong mọi
trường hợp, một cấu xạ đối hạt nhân là một toàn cấu nghiêm ngặt. Cũng lưu ý
rằng một toàn cấu nghiêm ngặt đồng thời là một đơn cấu thì là một đẳng cấu.
Ta để bạn đọc phát triển khái niệm đối ngẫu về \emph{đơn cấu nghiêm ngặt}.

Để làm rõ quan hệ giữa khái niệm cấu xạ $\mathcal F$-hạ xuống và khái niệm
toàn cấu nghiêm ngặt, ta đưa thêm vào các định nghĩa sau:

\result{ĐỊNH NGHĨA}{2.3} Một cấu xạ $\alpha:S'\longrightarrow S$ được gọi là
một \emph{toàn cấu phổ dụng} (tương ứng, một \emph{toàn cấu nghiêm ngặt phổ
dụng}) nếu, với mọi $T$ trên $S$, tích thớ $T'=S'\times_ST$ tồn tại và phép
chiếu $T'\longrightarrow T$ là một toàn cấu (tương ứng, một toàn cấu nghiêm
ngặt).

Trong các phạm trù rất tốt (như phạm trù các tập hợp, phạm trù các bó tập hợp
trên một không gian tôpô, các phạm trù abel, v.v.), bốn khái niệm về tính toàn
cấu vừa được đưa vào trùng nhau; trái lại, chúng đều khác nhau trong một phạm
trù như phạm trù các tiền lược đồ, hoặc phạm trù các tiền lược đồ trên một tiền
lược đồ không rỗng $S$ đã cho, ngay cả khi chỉ xét các $S$-lược đồ hữu hạn trên
$S$.

\result{ĐỊNH NGHĨA}{2.4} Một cấu xạ $\alpha:S'\longrightarrow S$ được gọi là
một \emph{cấu xạ hạ xuống} (tương ứng, một \emph{cấu xạ hạ xuống nghiêm
ngặt}) nếu nó là một cấu xạ $\mathcal F$-hạ xuống (tương ứng, một cấu xạ
$\mathcal F$-hạ xuống nghiêm ngặt) (xem Định nghĩa 1.7), trong đó $\mathcal F$
là phạm trù phân thớ trên cơ sở $\mathcal C$ gồm các đối tượng của
$\mathcal C$ trên các đối tượng của $\mathcal C$ (số 1, Ví dụ 1).

\result{MỆNH ĐỀ}{2.1} Nếu trong $\mathcal C$ tồn tại các tích hữu hạn và các
tích thớ (hữu hạn), thì trong $\mathcal C$, các cấu xạ hạ xuống chính là các
toàn cấu nghiêm ngặt phổ dụng.

\result{b. VÍ DỤ}{} Cho $\mathcal C$ là phạm trù các tiền lược đồ. Cho
$S\in\mathcal C$, và cho $S'$ và $S''$ là hai tiền lược đồ hữu hạn trên $S$,
nghĩa là tương ứng với các bó đại số $\mathcal A'$, $\mathcal A''$ trên $S$
mà, xét như các bó môđun, là tựa nhất quán và thuộc kiểu hữu hạn (nghĩa là
coherent nếu $S$ là Noether địa phương). Cho $\alpha:S'\longrightarrow S$ là
cấu xạ cấu trúc của $S'$, và cho $\beta_1,\beta_2$ là hai $S$-cấu xạ từ $S''$
đến $S'$, được xác định bởi các đồng cấu đại số
$\mathcal A'\longrightarrow\mathcal A''$, vẫn được ký hiệu là
$\beta_1,\beta_2$. Dùng sự kiện rằng một cấu xạ hữu hạn là đóng (định lý
Cohen--Seidenberg thứ nhất), ta dễ dàng chứng minh rằng biểu đồ
\[
  \tag{+}
  S\xleftarrow{\alpha}S'
  \overset{\beta_1}{\underset{\beta_2}{\rightrightarrows}}S''
\]
trong $\mathcal C$ là khớp khi và chỉ khi biểu đồ các bó

\src{306}
\[
  \mathcal O_S=\mathcal A\longrightarrow\mathcal A'
  \overset{\beta_1}{\underset{\beta_2}{\rightrightarrows}}\mathcal A''
\]
trên $S$ là khớp. Đặc biệt, nếu $\alpha:S'\longrightarrow S$ là một cấu xạ
hữu hạn tương ứng với một bó đại số $\mathcal A'$ trên $S$, thì $\alpha$ là
một toàn cấu nghiêm ngặt khi và chỉ khi biểu đồ các bó
\[
  \mathcal O_S=\mathcal A\longrightarrow\mathcal A'
  \overset{p_1}{\underset{p_2}{\rightrightarrows}}
  \mathcal A'\otimes_{\mathcal A}\mathcal A'
\]
là khớp (nó là một toàn cấu khi và chỉ khi
$\mathcal A\longrightarrow\mathcal A'$ là đơn ánh). Nếu $S$ là affine với
vành $A$, và do đó $S'$ là affine với vành $A'$ hữu hạn trên $A$, thì
$S'\longrightarrow S$ là một toàn cấu nghiêm ngặt khi và chỉ khi
$A\longrightarrow A'$ là một đẳng cấu từ $A$ lên vành con của $A'$ gồm những
$x'\in A'$ sao cho
\[
  1_{A'}\otimes_Ax'-x'\otimes_A1_{A'}=0
\]
(nó là một toàn cấu khi và chỉ khi $A\longrightarrow A'$ là đơn ánh). Như ta
đã lưu ý, ngay cả khi $S$ là lược đồ của một vành địa phương Artin, một cấu xạ
hữu hạn $S'\longrightarrow S$ là toàn cấu không nhất thiết là một toàn cấu
nghiêm ngặt. Tuy nhiên, có thể chứng minh rằng, nếu $S$ là một tiền lược đồ
Noether, thì mọi cấu xạ hữu hạn $S'\longrightarrow S$ là toàn cấu đều là hợp
thành của một dãy hữu hạn các toàn cấu nghiêm ngặt, mỗi cấu xạ trong dãy cũng
hữu hạn. Điều này cũng cho thấy hợp thành của hai toàn cấu nghiêm ngặt không
nhất thiết là một toàn cấu nghiêm ngặt.

\result{c}{} Nếu $(+)$ là một biểu đồ khớp gồm các cấu xạ hữu hạn, thì với mọi
cấu xạ phẳng $T\longrightarrow S$ của các tiền lược đồ, biểu đồ thu được từ
$(+)$ qua phép đổi cơ sở $T\longrightarrow S$ vẫn khớp. Suy ra rằng, nếu $X,Y$
là hai $S$-tiền lược đồ, với $X$ phẳng trên $S$, thì biểu đồ các ánh xạ giữa
tập hợp sau đây (trong đó $X',Y'$ là các ảnh ngược của $X,Y$ trên $S'$, còn
$X'',Y''$ là các ảnh ngược của chúng trên $S''$)
\[
  \operatorname{Hom}_S(X,Y)\longrightarrow
  \operatorname{Hom}_{S'}(X',Y')
  \rightrightarrows
  \operatorname{Hom}_{S''}(X'',Y'')
\]
là khớp. Đặc biệt, nếu $\mathcal F$ là phạm trù phân thớ trên cơ sở phạm trù
$\mathcal C$ các tiền lược đồ, sao cho với $X\in\mathcal C$, $\mathcal F_X$
là phạm trù các $X$-tiền lược đồ phẳng, thì biểu đồ $(+)$ là
$\mathcal F$-khớp. (Kết quả này không còn đúng nếu không đặt giả thiết phẳng;
đặc biệt, một toàn cấu nghiêm ngặt hữu hạn không nhất thiết là một cấu xạ hạ
xuống.) Tương tự, ta thấy $(+)$ là $\mathcal F$-khớp nếu $\mathcal F$ là phạm
trù phân thớ mà $\mathcal F_X$ là phạm trù các bó tựa nhất quán và phẳng trên
tiền lược đồ $X$ (ở đây giả thiết

\src{307}
phẳng cũng là thiết yếu). Trong cả hai trường hợp, vấn đề một dữ liệu dán có
hữu hiệu hay không (và cụ thể hơn, một dữ liệu hạ xuống khi
$S''=S'\times_SS'$ có hữu hiệu hay không) trên một đối tượng phẳng trên $S'$
là tinh tế (và lời giải của nó trong nhiều trường hợp riêng là một trong các
chủ đề chính của những bài trình bày này). Diễn giả không biết liệu, với mọi
toàn cấu nghiêm ngặt hữu hạn $S'\longrightarrow S$, mọi dữ liệu hạ xuống trên
một bó tựa nhất quán phẳng trên $S'$ có hữu hiệu hay không (ngay cả khi giả sử
$S$ là phổ của một vành địa phương Artin và chỉ xét các bó tự do địa phương
hạng 1). Tổng quát hơn, cho $A$ là một vành và $A'$ là một $A$-đại số (mọi
vành đều giao hoán) sao cho biểu đồ các ánh xạ
\[
  A\longrightarrow A'\rightrightarrows A'\otimes_AA'
\]
là khớp; điều này cũng tương đương với việc cấu xạ $S'\longrightarrow S$ tương
ứng giữa các phổ của $A,A'$ là một cấu xạ $\mathcal F$-hạ xuống, trong đó
$\mathcal F$ là phạm trù phân thớ các bó tựa nhất quán phẳng. Cho $M'$ là một
$A'$-môđun phẳng được trang bị một dữ liệu hạ xuống về $A$, nghĩa là một đẳng
cấu
\[
  \varphi:M'\otimes_AA'\xrightarrow{\sim}A'\otimes_AM'
\]
của các $A'\otimes_AA'$-môđun, thỏa mãn các điều kiện (i) và (ii) của 1.c (ta
để bạn đọc viết tường minh các điều kiện ấy bằng ngôn ngữ môđun). Dữ liệu này
có hữu hiệu hay không (đối với phạm trù phân thớ các bó tựa nhất quán phẳng)?
Cho $M$ là tập con của $M'$ gồm những $x'\in M'$ sao cho
\[
  \varphi(x'\otimes_A1_{A'})=1_{A'}\otimes_Ax',
\]
đó là một $A$-môđun con của $M'$. Phép nhúng chính tắc
$M\longrightarrow M'$ xác định một đồng cấu $A'$-môđun
$M\otimes_AA'\longrightarrow M'$. Khi đó tính hữu hiệu của $\varphi$ có nghĩa
là: $M$ là một $A$-môđun phẳng và đồng cấu trên là một đẳng cấu.

\result{NHẬN XÉT}{} Trong các lập luận trên, ta không đặt giả thiết phẳng nào
lên các cấu xạ của biểu đồ $(+)$; vì thế, để có một kỹ thuật hạ xuống, ta buộc
phải đặt các giả thiết phẳng lên những đối tượng trên $S,S'$ đang xét. Trong
§2, ta sẽ đặt giả thiết phẳng lên $\alpha:S'\longrightarrow S$, nhờ đó có được
một kỹ thuật hạ xuống cho những đối tượng trên $S,S'$ không còn phải chịu bất
kỳ điều kiện phẳng nào. Trong mọi trường hợp đều có một giả thiết phẳng tham
gia. Đây là một trong những lý do chính khiến khái niệm tính phẳng quan trọng
trong hình học đại số (vai trò này không thể lộ ra chừng nào người ta còn chỉ
xét các trường cơ sở, bởi trên đó, quả thực, mọi thứ đều phẳng!).

\src{308}
\subsection*{3. Ứng dụng vào các cấu xạ étale}
\addcontentsline{toc}{subsection}{3. Ứng dụng vào các cấu xạ étale}

Cho $A$ là một vành địa phương và $B$ là một đại số địa phương trên $A$ mà
iđêan cực đại của nó cảm sinh iđêan cực đại của $A$. Ta nói $B$ là
\emph{étale} trên $A$ (thay cho ``không phân nhánh'', cách gọi được dùng ở nơi
khác) nếu nó thỏa mãn các điều kiện sau:
\begin{enumerate}
\item[(i)] $B$ phẳng trên $A$;
\item[(ii)] $B/\mathfrak mB$ là một mở rộng hữu hạn khả phân của
$A/\mathfrak m=k$ (trong đó $\mathfrak m$ ký hiệu iđêan cực đại của $A$).
\end{enumerate}
Khi $A$ và $B$ là Noether và $k$ đóng đại số, điều này có nghĩa là đồng cấu
$\widehat A\longrightarrow\widehat B$ giữa các vành đầy đủ hóa, mở rộng
$A\longrightarrow B$, là một đẳng cấu. Một cấu xạ thuộc kiểu hữu hạn
$f:T\longrightarrow S$ được gọi là \emph{étale tại $x\in T$}, hay cũng nói
$T$ \emph{étale trên $S$ tại $x$}, nếu $\mathcal O_x$ étale trên
$\mathcal O_{f(x)}$; còn $f$ được gọi là \emph{étale}, hay $f$ được gọi là một
\emph{cấu xạ étale}, hay $T$ được gọi là \emph{étale trên $S$}, nếu $f$ étale tại mọi
$x\in T$. Hơn nữa, lưu ý rằng nếu $S$ là Noether địa phương, thì tập các điểm của
$T$ tại đó $f$ étale là một tập mở; việc dùng ``định lý chính'' của Zariski cho
phép mô tả chính xác cấu trúc của $T/S$ trong một lân cận của một điểm như thế
(bằng một phương trình thuộc kiểu quen thuộc).

Khi $S$ là một lược đồ thuộc kiểu hữu hạn trên trường các số phức, nó tương ứng
với một không gian giải tích $\overline S$ theo nghĩa của Serre [5], ngoại trừ
việc $\overline S$ có thể chứa các phần tử lũy linh trong bó cấu trúc của nó;
điều này không làm thay đổi gì cốt yếu trong [5]. Khi đó dễ thấy $f$ étale khi
và chỉ khi $\overline f:\overline T\longrightarrow\overline S$ étale, nghĩa là
mọi điểm của $\overline T$ có một lân cận mà trên đó $\overline f$ cảm sinh một
đẳng cấu lên một tập mở của $\overline S$. Đặc biệt, với mỗi \emph{phủ étale}
$T$ của $S$ (nghĩa là một cấu xạ hữu hạn étale $f:T\longrightarrow S$), ta có
một phủ étale $\overline T$ của $\overline S$, và phủ này liên thông khi và chỉ
khi $T$ liên thông [5]. Cũng dễ thấy rằng, nếu $T,T'$ là hai lược đồ
étale trên $S$, thì ánh xạ tự nhiên
\[
  \operatorname{Hom}_S(T,T')\longrightarrow
  \operatorname{Hom}_{\overline S}(\overline T,\overline{T'})
\]
là song ánh; nghĩa là hàm tử $T\longmapsto\overline T$ từ phạm trù các lược đồ
étale trên $S$ đến phạm trù các không gian giải tích étale trên $\overline S$
là ``đầy đủ và trung thành'', nên xác định một tương đương giữa phạm trù thứ
nhất và một phạm trù con của phạm trù thứ hai. Một định lý của
Grauert--Remmert [2] suy ra rằng, nếu $S$ chuẩn tắc, ta thu được một tương đương
giữa phạm trù các \emph{phủ étale} của $S$ và phạm trù các phủ étale (hữu hạn)
của $\overline S$; nghĩa là mọi phủ étale của $\overline S$ đều
$\overline S$-đẳng cấu với một $\overline T$, trong đó $T$ là một phủ étale của
$S$. Ta chứng minh rằng \emph{định lý Grauert--Remmert vẫn đúng mà không cần
giả thiết $S$ chuẩn tắc}. Trước hết, cho $S'\longrightarrow S$ là một toàn cấu
nghiêm ngặt hữu hạn; giả sử

\src{309}
định lý đã được chứng minh cho $S'$, và ta chứng minh nó đúng cho $S$. Thật vậy,
cho $\mathfrak T$ là một phủ étale của $\overline S$, và xét nghịch ảnh
$\mathfrak T'$ của nó trên $\overline{S'}$. Nghịch ảnh này tương ứng với một bó
giải tích coherent các đại số $\mathcal A'$ trên $\overline{S'}$, là nghịch ảnh
của bó các đại số $\mathcal A$ trên $\overline S$ xác định $\mathfrak T$. Theo giả
thiết, trên $\overline{S'}$, $\mathfrak T'$ đến từ một phủ étale $T'$ của $S'$,
nghĩa là $\mathcal A'$ đến từ một bó coherent các đại số $A'$ trên $S'$. Mặt
khác, $\mathcal A'$ được trang bị một dữ liệu hạ xuống chính tắc đối với
$\overline{S'}\longrightarrow\overline S$, nghĩa là một đẳng cấu giữa hai
nghịch ảnh của nó trên
\[
  \overline{S'}\times_{\overline S}\overline{S'}
  =\overline{(S'\times_SS')},
\]
(thỏa mãn các điều kiện (i), (ii)); theo những điều đã nói, đẳng cấu này đến từ
một đẳng cấu giữa các bó đại số tương ứng, nghĩa là từ một dữ liệu hạ xuống
trên $A'$ đối với $S'\longrightarrow S$. Dễ kiểm tra dữ liệu sau là hữu hiệu
(vì dữ liệu trên $\mathcal A'$ là hữu hiệu, và tính hữu hiệu của một dữ liệu hạ
xuống, như đã được giải thích tường minh ở mục trước, có thể được nhận biết địa
phương trên các \emph{đầy đủ hóa} của những môđun tham gia). Do đó ta thu
được một bó coherent các đại số $A$ trên $S$, xác định một phủ $T$ của $S$;
đây chính là phủ cần tìm. Kết quả trên hiển nhiên vẫn đúng nếu
$S'\longrightarrow S$ chỉ là hợp thành của một số hữu hạn các toàn cấu nghiêm
ngặt hữu hạn, nghĩa là một toàn cấu hữu hạn bất kỳ (theo kết quả phân tích thành
nhân tử được nêu trong mục 2). Suy ra định lý Grauert--Remmert vẫn đúng khi $S$
là một lược đồ \emph{rút gọn}, nghĩa là $\mathcal O_S$ không có phần tử lũy
linh, như ta thấy bằng cách đưa vào phép chuẩn hóa $S'$ của nó. Từ đó dễ dàng
chuyển sang trường hợp tổng quát.

Một chứng minh hoàn toàn tương tự, lại dùng kết quả phân tích thành nhân tử cho
các toàn cấu nghiêm ngặt hữu hạn và bản chất ``hình thức'' của tính hữu hiệu của
các dữ liệu hạ xuống, cho phép chứng minh kết quả sau: cho $S$ là một tiền lược
đồ Noether địa phương và $S'\longrightarrow S$ là một cấu xạ \emph{hữu hạn,
toàn ánh và thuần bất khả phân} (hay, tương đương, một cấu xạ thuộc kiểu hữu hạn
sao cho với mọi $T$ trên $S$, cấu xạ
$T'=S'\times_ST\longrightarrow T$ là một đồng phôi; điều này cũng được diễn đạt
bằng cách nói rằng $S'\longrightarrow S$ là một \emph{đồng phôi phổ dụng}). Với
mỗi $T$ étale trên $S$, xét nghịch ảnh $T'=T\times_SS'$ của nó, là étale trên
$S'$. Khi đó hàm tử $T\longmapsto T'$ là một \emph{tương đương phạm trù từ
phạm trù các tiền lược đồ $T$ étale trên $S$ đến phạm trù các tiền lược đồ $T'$
étale trên $S'$}. (Ta dùng tính song ánh của
\[
  \operatorname{Hom}_S(T_1,T_2)\longrightarrow
  \operatorname{Hom}_{S'}(T'_1,T'_2)
\]
đối với hai tiền lược đồ $T_1,T_2$ étale trên $S$, một sự kiện dễ kiểm tra trực
tiếp, và sự kiện rằng định lý vừa nêu đúng nếu
$S'=(S,\mathcal O_S/\mathcal J)$,

\src{310}
trong đó $\mathcal J$ là một bó coherent các iđêan lũy linh của $\mathcal O_S$
([4], Bổ đề 6)). Hơn nữa, lưu ý rằng ở đây ta không giả sử các $T$ đang xét hữu
hạn trên $S$. Đặc biệt, kết quả này suy ra cấu xạ $S'\longrightarrow S$ cảm
sinh một đẳng cấu từ nhóm cơ bản $[\,]$ của $S'$ lên nhóm cơ bản của $S$
(``\emph{tính bất biến tôpô của nhóm cơ bản của một tiền lược đồ}'').

\subsection*{4. Quan hệ với đối đồng điều bậc 1}
\addcontentsline{toc}{subsection}{4. Quan hệ với đối đồng điều bậc 1}

\medskip\noindent
a. Cho $\mathcal C$ là một phạm trù trong đó tích của hai đối tượng bất kỳ luôn
tồn tại, và cho $T\in\mathcal C$. Với mỗi tập hữu hạn
$I\ne\varnothing$, ta có thể xét $T^I$; khi $I$ biến thiên, ta thu được một hàm
tử hiệp biến từ phạm trù các tập hữu hạn khác rỗng đến $\mathcal C$, nghĩa là
một cái có thể gọi là \emph{đối tượng đơn hình} của $\mathcal C$, ký hiệu
$K_T$. Đối tượng này phụ thuộc hiệp biến vào $T$; hơn nữa, nếu $u,v$ là hai
cấu xạ $T\longrightarrow T'$, thì các cấu xạ tương ứng
$K_T\longrightarrow K_{T'}$ đồng luân với nhau. Ta nói $T$
\emph{trội hơn} $T'$ nếu $\operatorname{Hom}(T,T')\ne\varnothing$; đây là một
quan hệ tiền thứ tự lọc tăng trên $\mathcal C$. Từ những điều trên suy ra rằng,
nếu $T$ trội hơn $T'$, thì có một lớp chính tắc (sai khác một đồng luân) các
đồng cấu giữa những đối tượng đơn hình $K_T\longrightarrow K_{T'}$; đặc biệt,
nếu $K_T$ và $K_{T'}$ sao cho mỗi đối tượng trội hơn đối tượng kia, thì $K_T$
và $K_{T'}$ tương đương đồng luân. Bây giờ cho $F$ là một hàm tử (phản biến,
để cố định ý tưởng) từ $\mathcal C$ đến một phạm trù Abel $\mathcal C'$; khi đó
\[
  C^*(T,F)=F(K_T)
\]
là một đối tượng đối đơn hình của $\mathcal C'$, và do đó xác định theo cách
quen thuộc một phức (đối dây chuyền) trong $\mathcal C'$, mà ta có thể lấy đối
đồng điều:
\[
  H^*(T,F)=H^*(C^*(T,F))=H^*(F(K_T))
\]
(Ta có thể đặt một $\mathcal C$ làm chỉ số dưới của $H^*$ nếu có khả năng nhầm
lẫn.) Đây là một hàm tử đối đồng điều theo $F$; tính hiệp biến hoặc phản biến
của nó khi $T$ biến thiên suy ra từ những điều đã nói về $K_T$. Chính xác hơn,
với $F$ cố định và $T$ biến thiên trong $\mathcal C$ (được tiền thứ tự hóa bởi
quan hệ trội hơn), các $H^*(T,F)$ lập thành một hệ quy nạp các đối tượng phân
bậc của $\mathcal C'$; đặc biệt, nếu $T$ và $T'$ sao cho mỗi đối tượng trội hơn
đối tượng kia, thì $H^*(T,F)$ và $H^*(T',F)$ đẳng cấu chính tắc.

Giả sử các tích thớ tồn tại trong $\mathcal C$. Với $S\in\mathcal C$ cố định,
ta có thể áp dụng những điều trên cho phạm trù $\mathcal C_S$ gồm các đối tượng
của $\mathcal C$ trên $S$; ta sẽ viết $C^*(T/S,F)$ và $H^*(T/S,F)$ thay cho
$C^*(T,F)$ và $H^*(T,F)$ khi muốn nói rõ rằng ta đang làm việc trong phạm trù
$\mathcal C_S$; như vậy,

\src{311}
$C^*(T/S,F)$ là một phức đối dây chuyền trong $\mathcal C'$ mà số hạng ở bậc
$n$ bằng
\[
  F(T\times_ST\times_S\cdots\times_ST)
\]
(trong đó biểu thức trong ngoặc có $n+1$ thừa số).

Như thường lệ, lưu ý rằng ta có thể định nghĩa $H^0(T/S,F)$ mà không giả sử
phạm trù $\mathcal C'$ là Abel: đó là hạt nhân (Định nghĩa 2.1), nếu tồn tại,
của cặp cấu xạ $F(p_i)$ ($i=1,2$)
\[
  F(T)\rightrightarrows F(T\times_ST)
\]
tương ứng với hai phép chiếu
$p_1,p_2:T\times_ST\rightrightarrows T$. Đặc biệt, ta có một cấu xạ tự nhiên
(gọi là \emph{cấu xạ tăng cường})
\[
  F(S)\longrightarrow H^0(T/S,F)
\]
là một đẳng cấu trong các trường hợp thuận lợi (đặc biệt, nếu
$T\longrightarrow S$ là một toàn cấu nghiêm ngặt và $F$ “khớp trái”).
Tương tự, khi $F$ nhận giá trị trong phạm trù các nhóm trong một phạm trù
$\mathcal C''$, ta cũng có thể định nghĩa $H^1(T/S,F)$. Khi $\mathcal C''$ là
phạm trù các tập hợp (nghĩa là $F$ nhận giá trị trong phạm trù các nhóm thông
thường, không nhất thiết giao hoán), $H^1(T/S,F)$ là tập thương của nhóm con
$Z^1(T/S,F)$ của $C^1(T/S,F)=F(T\times_ST)$ gồm các phần tử $g$ thỏa mãn
\[
  F(p_{31})(g)=F(p_{32})(g)\,F(p_{21})(g)
\]
theo nhóm toán tử $F(T)$, tác động lên $C^1(T/S,F)$ và, đặc biệt, lên tập con
$Z^1(T/S,F)$ bởi
\[
  \rho(g')\cdot g=F(p_2)(g')\,g\,F(p_1)(g')^{-1}.
\]

\medskip\noindent
b. Chẳng hạn, cho $\mathcal F$ là một phạm trù phân thớ trên cơ sở
$\mathcal C$. Cho $\xi,\eta\in\mathcal F_S$ và, với mỗi $S'$ trên $S$, đặt
\[
  F_{\xi,\eta}(S')=
  \operatorname{Hom}(\xi\times_SS',\eta\times_SS').
\]
Như vậy $F_{\xi,\eta}$ là một hàm tử phản biến từ $\mathcal C_S$ đến phạm trù
các tập hợp. Với ký hiệu đó, \emph{nói rằng cấu xạ tăng cường}
\[
  F_{\xi,\eta}(S)\longrightarrow H^0(S'/S,F_{\xi,\eta})
\]
\emph{là một đẳng cấu với mọi cặp phần tử
$\xi,\eta\in\mathcal F_S$ có nghĩa là
$\alpha:S'\longrightarrow S$ là một cấu xạ $\mathcal F$-hạ xuống}
(Định nghĩa 1.7).

\medskip\noindent
c. Tương tự, với $\xi\in\mathcal F_S$ và mỗi đối tượng $S'$ của $\mathcal C$
trên $S$, đặt
\[
  G_\xi(S')=\operatorname{Aut}(\xi\times_SS') ;
\]
như vậy ta đã định nghĩa một hàm tử phản biến $G_\xi$ từ $\mathcal C_S$ đến
phạm trù các

\src{312}
nhóm. Khi đó ta thấy $Z^1(S'/S,G_\xi)$ \emph{được đồng nhất chính tắc với tập
các dữ liệu hạ xuống trên $\xi'=\xi\times_SS'$ đối với
$S'\longrightarrow S$} (Định nghĩa 1.6), còn $H^1(S'/S,G_\xi)$ \emph{được
đồng nhất với tập các lớp, sai khác một đẳng cấu, của những đối tượng trong
$\mathcal F_{S'}$ được trang bị một dữ liệu hạ xuống đối với
$\alpha:S'\longrightarrow S$ và, xét như các đối tượng của
$\mathcal F_{S'}$, đẳng cấu với $\xi'=\xi\times_SS'$}. Do đó, nếu
$\alpha:S'\longrightarrow S$ là một cấu xạ $\mathcal F$-hạ xuống (xem (b)),
thì $H^1(S'/S,G_\xi)$ \emph{chứa như một tập con tập các lớp, sai khác một
đẳng cấu, của những đối tượng $\eta$ trong $\mathcal F_S$ sao cho
$\eta\times_SS'$ đẳng cấu trong $\mathcal F_{S'}$ với $\xi\times_SS'$; và phép
nhúng này là ánh xạ đồng nhất khi và chỉ khi mọi dữ liệu hạ xuống trên
$\xi'=\xi\times_SS'$ đối với $\alpha:S'\longrightarrow S$ đều hữu hiệu}.
(Điều này đặc biệt đúng nếu $\alpha:S'\longrightarrow S$ là một cấu xạ
$\mathcal F$-hạ xuống nghiêm ngặt.)

\result{NHẬN XÉT}{} Các phức đối dây chuyền dạng $C^*(T/S,F)$ chứa, như những
trường hợp riêng, phần lớn các phức chuẩn quen thuộc (đối đồng điều Čech, đối
đồng điều nhóm, v.v.), và giữ một vai trò quan trọng trong hình học đại số,
đặc biệt trong “đối đồng điều Weil” của các tiền lược đồ.

\medskip\noindent
d. \result{VÍ DỤ}{1} Cho $S'$ là một đối tượng trên $S\in\mathcal C$, và cho
$\Gamma$ là một nhóm tự đẳng cấu của $S'$ sao cho $S'$ “là chính trên $S$
theo nghĩa hình thức, với nhóm $\Gamma$”, nghĩa là cấu xạ tự nhiên
\[
  \Gamma\times S'\longrightarrow S'\times_SS',
\]
trong đó $\Gamma\times S'$ chỉ tổng trực tiếp của $\Gamma$ bản sao của $S'$,
là một đẳng cấu. (Ta giả sử các tổng trực tiếp xuất hiện
ở đây tồn tại trong $\mathcal C$.) Cho $F$ là một hàm tử phản biến từ
$\mathcal C$ đến phạm trù các nhóm Abel. Khi đó $C^*(S'/S,F)$ đẳng cấu chính
tắc với nhóm đơn hình của các đối dây chuyền thuần nhất chuẩn
$C^*(\Gamma,F(S'))$; do đó $H^*(S'/S,F)$ đẳng cấu chính tắc với
$H^*(\Gamma,F(S'))$.

\medskip\noindent
e. \result{VÍ DỤ}{2} Cho $\mathcal C$ là phạm trù các tiền lược đồ. Ký hiệu
$G_a$ (“nhóm cộng”) là hàm tử phản biến từ $\mathcal C$ đến phạm trù các
nhóm Abel, được định nghĩa bởi
\[
  G_a(X)=H^0(X,\mathcal O_X).
\]
Tương tự, ta định nghĩa hàm tử $G_m$ (“nhóm nhân”) bởi
\[
  G_m(X)=H^0(X,\mathcal O_X)^*
\]
(nhóm các phần tử khả nghịch của vành $H^0(X,\mathcal O_X)$), và tổng quát hơn,
hàm tử $\operatorname{Gl}(n)$ (“nhóm tuyến tính cấp $n$”) bởi

\src{313}
\[
  \operatorname{Gl}(n)(X)=
  \operatorname{Gl}\bigl(n,H^0(X,\mathcal O_X)\bigr),
\]
là một hàm tử từ $\mathcal C$ đến phạm trù các nhóm (không nhất thiết giao
hoán nếu $n>1$; với $n=1$ ta thu lại $G_m$). Hơn nữa,
$\operatorname{Gl}(n)$ có thể được diễn giải như một hàm tử các tự đẳng cấu (xem
(c)) bằng cách xét phạm trù phân thớ $\mathcal F$ trên cơ sở $\mathcal C$ sao
cho, với $X\in\mathcal C$, $\mathcal F_X$ là phạm trù các bó tự do địa phương
trên $X$: quả vậy,
\[
  \operatorname{Gl}(n)(X)=
  \operatorname{Aut}_{\mathcal F_X}(\mathcal O_X^n).
\]
Theo (b), suy ra rằng nếu $\alpha:S'\longrightarrow S$ là một cấu xạ
$\mathcal F$-hạ xuống (xem đoạn 2.c), thì
$H^1(S'/S,\operatorname{Gl}(n))$ chứa tập các lớp, sai khác một đẳng cấu, của
những bó tự do địa phương trên $S$ mà nghịch ảnh trên $S'$ đẳng cấu với
$\mathcal O_{S'}^n$; và phép nhúng này là một đẳng thức khi và chỉ khi mọi dữ
liệu hạ xuống trên $\mathcal O_{S'}^n$ (đối với
$\alpha:S'\longrightarrow S$) đều hữu hiệu. Khi $S$ là phổ của một vành địa
phương, điều đó có nghĩa là
$H^1(S'/S,\operatorname{Gl}(n))=(e)$, vì mọi bó tự do địa phương trên $S$ khi
ấy đều tầm thường.

Lưu ý tính tương đương của các điều kiện sau đây đối với một cấu xạ
$\alpha:S'\longrightarrow S$:
\begin{enumerate}
\item[(i)] đồng cấu tăng cường
\[
  H^0(S,\mathcal O_S)=G_a(S)\longrightarrow H^0(S'/S,G_a)
\]
là một đẳng cấu;
\item[(ii)] $S'\longrightarrow S$ là một cấu xạ $\mathcal F$-hạ xuống
(trong đó $\mathcal F$ là phạm trù phân thớ trên cơ sở $\mathcal C$ đã xét ở
trên).
\end{enumerate}
Nếu $S'\longrightarrow S$ là hữu hạn, các điều kiện này còn tương đương với:
\begin{enumerate}
\item[(iii)] $S'\longrightarrow S$ là một toàn cấu nghiêm ngặt
(xem đoạn 2.c).
\end{enumerate}

Bây giờ giả sử $S=\operatorname{Spec}(A)$ và
$S'=\operatorname{Spec}(A')$. Khi đó
\[
  C^n(S'/S,G_a)=C^n(A'/A,G_a)=\bigotimes_A^{n+1}A',
\]
trong đó toán tử đối biên
$C^n(A'/A,G_a)\longrightarrow C^{n+1}(A'/A,G_a)$ là tổng luân phiên của các
toán tử mặt
\[
  \partial_i(x_0\otimes x_1\otimes\cdots\otimes x_n)
  =x_0\otimes\cdots\otimes x_{i-1}\otimes1_{A'}\otimes
   x_i\otimes\cdots\otimes x_n.
\]
Tương tự,
\[
  C^n(S'/S,G_m)=C^n(A'/A,G_m)
  \simeq\left(\bigotimes_A^{n+1}A'\right)^*,
\]
các phép toán đơn hình trong $C^*(A'/A,G_m)$ được cảm sinh bởi những phép toán
trong $C^*(S'/S,G_a)$. Các phép toán đơn hình trong
$C^*(A'/A,\operatorname{Gl}(n))$ cũng được viết tường minh theo cách ấy.
\emph{Trong mọi trường hợp mà báo cáo viên biết, ta có}
\[
  H^i(A'/A,G_a)=0\qquad\text{với $i>0$},
\]
\emph{và nếu $A$ là địa phương, ta có}
\[
  H^1(A'/A,G_m)=0
  \quad\text{và tổng quát hơn}\quad
  H^1(A'/A,\operatorname{Gl}(n))=(e)
\]
(khi $S'\longrightarrow S$ là một cấu xạ $\mathcal F$-hạ xuống, nghĩa là biểu
đồ
$A\longrightarrow A'\rightrightarrows A'\otimes_AA'$ khớp; so sánh với đoạn
2.c). Lưu ý rằng \emph{“định lý 90” của Hilbert chẳng qua là hệ thức}

\src{314}
\[
  H^1(S'/S,G_m)=0
\]
\emph{khi $A$ là một trường và $A'$ là một mở rộng Galois hữu hạn của trường ấy
(xem Ví dụ 1), và còn có thể được phát biểu rằng, trong trường hợp đang xét,
$S'\longrightarrow S$ là một cấu xạ hạ xuống nghiêm ngặt đối với phạm trù phân
thớ các bó tự do địa phương hạng 1.} Chính dưới dạng cuối này mà định lý Hilbert
nên được tổng quát hóa, bằng cách thay đổi cả các giả thiết về cấu xạ
$S'\longrightarrow S$ lẫn các giả thiết về những bó tựa nhất quán đang xét.

Cuối cùng, lưu ý tính tương đương của các tính chất sau khi $A$ là một vành
Artin địa phương có iđêan cực đại $\mathfrak m$ và $A'$ là một $A$-đại số
(trong đó, với mỗi số nguyên $k>0$, $A_k$ (tương ứng $A'_k$) ký hiệu vành
$A/\mathfrak m^{k+1}$ (tương ứng $A'/\mathfrak m^{k+1}A'$)):
\begin{enumerate}
\item[(i)] $H^1(A'_k/A_k,G_a)=0$ với mọi $k$.
\item[(ii)] $H^1(A'_k/A_k,G_m)=0$ với mọi $k$.
\item[(iii)] $H^1(A'_k/A_k,\operatorname{Gl}(n))=(e)$ với mọi $k$ và mọi
$n$.
\end{enumerate}
Nếu $S'\longrightarrow S$ là một toàn cấu nghiêm ngặt, thì các điều kiện trên
thậm chí còn kéo theo rằng đó là một cấu xạ hạ xuống nghiêm ngặt đối với các
môđun tự do trên $A'$ (thuộc kiểu hữu hạn hoặc không).

\result{NHẬN XÉT}{} Định nghĩa các nhóm $H^i(S'/S,G_m)$, khi $S,S'$ là phổ của
các trường $A,A'$, là do Amitsur. Nhóm $H^2(S'/S,G_m)$ đặc biệt đáng quan tâm
như một biến thể “toàn cục” của nhóm Brauer; về biến thể này có thể tham khảo
[1], Chương VII.

\fsection{B}{Hạ xuống theo các cấu xạ phẳng trung thành}

\subsection*{1. Phát biểu các định lý hạ xuống}
\addcontentsline{toc}{subsection}{1. Phát biểu các định lý hạ xuống}

\result{ĐỊNH NGHĨA}{1.1} Một cấu xạ
$\alpha:S'\longrightarrow S$ của các tiền lược đồ được gọi là \emph{phẳng} nếu,
với mọi $x'\in S'$, $\mathcal O_{x'}$ là một môđun phẳng trên vành
$\mathcal O_{\alpha(x')}$ (nghĩa là
$\mathcal O_{x'}\otimes_{\mathcal O_{\alpha(x')}}M$ là một hàm tử khớp theo
$\mathcal O_{\alpha(x')}$-môđun $M$). Một cấu xạ được gọi là \emph{phẳng trung
thành} nếu nó phẳng và toàn ánh.

Chẳng hạn, nếu $S=\operatorname{Spec}(A)$ và
$S'=\operatorname{Spec}(A')$, thì $S'$ phẳng trên $S$ khi và chỉ khi $A'$ là
một $A$-môđun phẳng, và $S'$ phẳng trung thành trên $S$ khi và chỉ khi $A'$ là
một $A$-môđun phẳng trung thành (nghĩa là hàm tử $A'\otimes_A M$ theo
$A$-môđun $M$ là khớp và trung thành); theo thuật ngữ của Serre [5], điều này
cũng có nghĩa là cặp $(A,A')$ phẳng. Nếu $S'$ phẳng trung thành trên $S$, thì
hàm tử ảnh ngược trên các bó tựa nhất quán trên $S$ là khớp và trung thành;
nói cách khác, để một dãy các đồng cấu của các bó

\src{315}
tựa nhất quán trên $S$ là khớp, điều kiện cần và đủ là ảnh ngược của nó trên
$S'$ là khớp (đặc biệt, để một đồng cấu của các bó tựa nhất quán trên $S$ là
đơn cấu, tương ứng là toàn cấu, tương ứng là đẳng cấu, điều kiện cần và đủ là
ảnh ngược của nó trên $S'$ có tính chất ấy). Tính chất này vẫn đúng khi hạn chế
trên một phần mở tùy ý của $S$ và, dưới dạng đó, nó đặc trưng các cấu xạ phẳng
trung thành.

\result{ĐỊNH NGHĨA}{1.2} Một cấu xạ $\alpha:S'\longrightarrow S$ được gọi là
\emph{quasi-compact} nếu ảnh ngược của mọi phần mở quasi-compact $U$ của $S$
là quasi-compact (nghĩa là hợp hữu hạn của các phần mở affine).

Hiển nhiên, chỉ cần kiểm tra tính chất này đối với các phần mở affine của $S$.
Chẳng hạn, một cấu xạ affine (nghĩa là cấu xạ sao cho ảnh ngược của một phần mở
affine là affine) là quasi-compact.

Lớp các cấu xạ phẳng, tương ứng phẳng trung thành, tương ứng quasi-compact, ổn
định qua phép hợp thành và qua ``đổi cơ sở'', và dĩ nhiên chứa các đẳng cấu.

\result{ĐỊNH LÝ}{1} Cho $\alpha:S'\longrightarrow S$ là một cấu xạ phẳng trung
thành và quasi-compact của các tiền lược đồ. Khi đó $\alpha$ là một \emph{cấu
xạ hạ xuống nghiêm ngặt} (A, Định nghĩa 1.7) đối với phạm trù phân thớ
$\mathcal F$ của các bó tựa nhất quán (A, Mục 1, Ví dụ 2).

Phát biểu này có hai ý:
\begin{enumerate}
\item[(i)] Nếu $F$ và $G$ là hai bó tựa nhất quán trên $S$, còn $F'$ và $G'$
là các ảnh ngược của chúng trên $S'$, thì đồng cấu tự nhiên
\[
  \operatorname{Hom}(F,G)\longrightarrow\operatorname{Hom}(F',G')
\]
là một song ánh từ vế thứ nhất lên nhóm con của vế thứ hai gồm các đồng cấu
$F'\longrightarrow G'$ tương thích với các dữ liệu hạ xuống chính tắc trên các
bó ấy, nghĩa là những đồng cấu mà ảnh ngược theo hai phép chiếu từ
$S''=S'\times_SS'$ lên $S'$ cho cùng một đồng cấu
$F''\longrightarrow G''$.
\item[(ii)] Mọi bó tựa nhất quán $F'$ trên $S'$, được trang bị một dữ liệu hạ
xuống đối với cấu xạ $\alpha:S'\longrightarrow S$ (A, Định nghĩa 1.6), đều
đẳng cấu (cùng với dữ liệu đó) với ảnh ngược của một bó tựa nhất quán $F$ trên
$S$.
\end{enumerate}
Lấy $F=\mathcal O_S$ trong (i), ta được:

\result{HỆ QUẢ}{1} Cho $G$ là một bó tựa nhất quán trên $S$; gọi $G'$ và $G''$
là các ảnh ngược của nó trên $S'$ và trên $S''=S'\times_SS'$, và gọi $p_1,p_2$
là hai phép chiếu từ $S''$ lên $S'$. Khi đó biểu đồ các ánh xạ tập hợp sau

\src{316}
\[
  \Gamma(G)\xrightarrow{\alpha^*}\Gamma(G')
  \underset{p_2^*}{\overset{p_1^*}{\rightrightarrows}}\Gamma(G'')
\]
là khớp (A, Định nghĩa 1.(a)).

Mặt khác, kết hợp (i), (ii) và Định nghĩa 1.1 ta được

\result{HỆ QUẢ}{2} Cho $G$ như trong Hệ quả 1. Khi đó có một tương ứng song ánh
giữa các bó con tựa nhất quán của $G$ và các bó con tựa nhất quán của $G'$ mà
ảnh ngược của chúng trên $S''$ theo hai phép chiếu $p_1,p_2$ cho cùng một bó
con của $G''$.

Dĩ nhiên có một phát biểu tương đương bằng ngôn ngữ các bó thương. Như đã thấy
ở A, Mục 4, (e), Định lý 1 phải được xem là một sự khái quát hóa ``Định lý 90''
của Hilbert, và nó bao hàm như những trường hợp riêng nhiều cách phát biểu bằng
đối đồng điều bậc 1. Để chứng minh, ta dễ dàng quy về trường hợp
$S=\operatorname{Spec}(A)$, $S'=\operatorname{Spec}(A')$; và đối với (i), ta
dễ dàng quy về việc chứng minh Hệ quả 1, nghĩa là tính khớp của biểu đồ
\[
  M=A\otimes_A M\longrightarrow A'\otimes_A M
  \rightrightarrows A'\otimes_AA'\otimes_A M
\]
đối với mọi $A$-môđun $M$, điều này suy ra từ bổ đề tổng quát hơn sau đây:

\result{BỔ ĐỀ}{1.1} Cho $A'$ là một $A$-đại số phẳng trung thành. Khi đó, với
mọi $A$-môđun $M$, phức được tăng cường bởi $M$
$C^*(A'/A,G_a)\otimes_A M$ (xem A, Mục 4, (e)) là một \emph{phân giải} của
$M$.

Chỉ cần chứng minh rằng phức tăng cường thu được từ phức trên bằng phép đổi cơ
sở từ $A$ sang $A'$ thỏa mãn các kết luận tương tự. Điều này dẫn đến việc kiểm
tra phát biểu sau khi thay $A$ bởi $A'$ và $A'$ bởi $A'\otimes_AA'$, do đó quy
ta về trường hợp có một đồng cấu $A$-đại số $A'\longrightarrow A$ (hay, bằng
ngôn ngữ hình học, trường hợp $S'$ trên $S$ có một tiết diện). Trong trường hợp
ấy, khẳng định suy ra từ các sự kiện tổng quát ở A, Mục 4, (a). Nhân tiện, ta
nêu hệ quả sau, khái quát hóa một phát biểu quen thuộc trong đối đồng điều
Galois (so sánh A, Mục 4, (e)):

\result{HỆ QUẢ}{} Nếu $A'$ phẳng trung thành trên $A$, thì
$H^0(A'/A,G_a)=A$ và
\[
  H^i(A'/A,G_a)=0\qquad\text{với $i\geq1$}.
\]

Để chứng minh phần (ii) của Định lý 1, ta tiến hành như đối với (i), quy về
trường hợp $S'$ trên $S$ có một tiết diện; khi đó khẳng định suy ra từ (i)
(xem A, Mục 1, (c)).

Hiển nhiên, có thể biến đổi tùy ý Định lý 1 và các hệ quả của nó bằng cách đưa
thêm nhiều cấu trúc khác nhau lên các bó (hoặc các hệ bó) tựa nhất quán đang
xét. Chẳng hạn, việc cho trên $S$ một bó tựa nhất quán các đại số giao hoán

\src{317}
``tương đương'' với việc cho trên $S'$ một bó cùng loại được trang bị dữ liệu
hạ xuống đối với $\alpha:S'\longrightarrow S$. Xét đến sự tương ứng hàm tử
giữa các bó tựa nhất quán như vậy trên $S$ và các tiền lược đồ affine trên $S$,
ta thu được khẳng định thứ hai của định lý sau:

\result{ĐỊNH LÝ}{2} Cho $\alpha:S'\longrightarrow S$ như trong Định lý 1. Khi
đó $\alpha$ là một \emph{cấu xạ hạ xuống} (nói chung không nghiêm ngặt)
(A, Định nghĩa 2.4), và là một \emph{cấu xạ hạ xuống nghiêm ngặt} đối với phạm
trù phân thớ gồm các lược đồ affine trên các tiền lược đồ (A, Định nghĩa 1.7).

Khẳng định thứ nhất của định lý có nghĩa như sau: cho $X,Y$ là hai tiền lược đồ
trên $S$; gọi $X',Y'$ là các ảnh ngược của chúng trên $S'$, và gọi $X'',Y''$
là các ảnh ngược của chúng trên $S''=S'\times_SS'$. Khi đó biểu đồ các ánh xạ
tự nhiên sau
\[
  \operatorname{Hom}_S(X,Y)\xrightarrow{\alpha^*}
  \operatorname{Hom}_{S'}(X',Y')
  \underset{p_2^*}{\overset{p_1^*}{\rightrightarrows}}
  \operatorname{Hom}_{S''}(X'',Y'')
\]
là khớp; nghĩa là $\alpha^*$ là một song ánh từ
$\operatorname{Hom}_S(X,Y)$ lên phần của $\operatorname{Hom}_{S'}(X',Y')$ gồm
các đồng cấu tương thích với các dữ liệu hạ xuống chính tắc trên $X',Y'$ (nghĩa
là những đồng cấu mà ảnh ngược theo hai phép chiếu từ $S''$ lên $S'$ bằng
nhau). Điều này dễ dàng suy ra từ Định lý 1, Hệ quả 1, khi hạn chế $Y$ là affine
trên $S$; trong trường hợp tổng quát, phải kết hợp Định lý 1 với kết quả sau:

\result{BỔ ĐỀ}{1.2} Cho $\alpha:S'\longrightarrow S$ là một cấu xạ phẳng trung
thành và quasi-compact. Khi đó $S$ đồng nhất với một \emph{không gian tôpô
thương} của $S'$; nghĩa là mọi phần $U$ của $S$ sao cho $\alpha^{-1}(U)$ mở thì
chính nó cũng mở.

Để hoàn tất Định lý 2, cần cho các tiêu chuẩn về tính hữu hiệu của một dữ liệu
hạ xuống trên một $S'$-tiền lược đồ $X'$ (khi không giả sử $X'$ affine trên
$S'$). Trước hết, lưu ý rằng \emph{một dữ liệu hạ xuống như vậy không nhất
thiết hữu hiệu}, ngay cả khi $S$ là phổ của một trường $k$, $S'$ là phổ của một
mở rộng bậc hai $k'$ của trường ấy, và $X'$ là một lược đồ đại số proper chiều 2 trên
$S'$ (như có thể thấy, theo Serre, bằng cách dùng mặt không xạ ảnh của Nagata).
\emph{Để một dữ liệu hạ xuống trên $X'/S'$ đối với
$\alpha:S'\longrightarrow S$ (phẳng trung thành và quasi-compact) là hữu hiệu,
điều kiện cần và đủ là $X'$ là hợp của các phần mở $X'_i$, affine trên $S'$, và
``ổn định'' dưới dữ liệu hạ xuống trên $X'$.} Điều này chắc chắn đúng (bất kể
$X'/S'$ và dữ liệu hạ xuống trên $X'$ là gì) nếu cấu xạ
$\alpha:S'\longrightarrow S$ là \emph{radicial} (nghĩa là đơn ánh và các mở
rộng trường thặng dư

\src{318}
cũng radicial). Cũng có thể chứng minh rằng điều ấy vẫn đúng nếu
$\alpha:S'\longrightarrow S$ là \emph{hữu hạn} và mọi phần hữu hạn của $X'$,
nằm trong một thớ của $X'$ trên $S$, đều nằm trong một phần mở của $X'$ affine
trên $S$ (đó là \emph{tiêu chuẩn Weil}). Đặc biệt, điều này đúng nếu $X'/S'$ là
quasi-xạ ảnh, và trong trường hợp ấy có thể chứng minh tiền lược đồ ``hạ xuống''
$X/S$ cũng quasi-xạ ảnh (và xạ ảnh nếu $X'/S'$ xạ ảnh). Tóm lại:

\result{ĐỊNH LÝ}{3} Cho $\alpha:S'\longrightarrow S$ là một cấu xạ phẳng trung
thành và quasi-compact của các tiền lược đồ. Nếu $\alpha$ là \emph{radicial},
thì nó là một \emph{cấu xạ hạ xuống nghiêm ngặt}. Nếu $\alpha$ hữu hạn, thì nó
là một cấu xạ hạ xuống nghiêm ngặt đối với phạm trù phân thớ gồm các tiền lược
đồ quasi-xạ ảnh (hoặc xạ ảnh) trên các tiền lược đồ.

\result{NHẬN XÉT}{} Tôi không biết trong khẳng định thứ hai ở trên, giả thiết
$\alpha$ là một cấu xạ \emph{hữu hạn} có thực sự cần thiết hay không; dù sao,
có thể kiểm tra một cách ``hình thức'' rằng giả thiết ấy có thể được thay bằng
giả thiết sau, có vẻ yếu hơn: \emph{với mọi điểm của $S$, tồn tại một lân cận
mở $U$, một $U'$ hữu hạn và phẳng trung thành trên $U$, cùng một $S$-cấu xạ từ
$U'$ vào $S'$.} Một trường hợp điển hình không nằm trong trường hợp trên là
$S=\operatorname{Spec}(A)$, $S'=\operatorname{Spec}(\widehat A)$, trong đó $A$
là một vành Noether địa phương và $\widehat A$ là phép đầy đủ hóa của nó; hoặc
trường hợp $S'$ quasi-hữu hạn trên $S$ (nghĩa là đẳng cấu địa phương với một
phần mở của một $S$-lược đồ hữu hạn) nhưng không hữu hạn. Trong hai trường hợp
ấy, người trình bày cũng không biết câu trả lời cho câu hỏi sau: cho $X$ là một
$S$-lược đồ sao cho $X'=X\times_SS'$ xạ ảnh trên $S'$; liệu $X$ có xạ ảnh trên
$S$ hay không?

\subsection*{2. Áp dụng vào việc hạ xuống một số tính chất của cấu xạ}
\addcontentsline{toc}{subsection}{2. Hạ xuống các tính chất của cấu xạ}

Cho $P$ là một lớp các cấu xạ của tiền lược đồ. Cho
$\alpha:S'\longrightarrow S$ là một cấu xạ của tiền lược đồ, cho
$f:X\longrightarrow Y$ là một cấu xạ của các $S$-tiền lược đồ, và cho
$f':X'\longrightarrow Y'$ là nghịch ảnh của $f$ qua $\alpha$. Khi đó ta có thể
hỏi liệu quan hệ ``$f'\in P$'' có kéo theo ``$f\in P$'' hay không. Trong nhiều
trường hợp quan trọng, câu trả lời là khẳng định nếu giả sử $\alpha$
\emph{phẳng trung thành và quasi-compact} (giả thiết sau đôi khi là thừa). Điều
này thấy được trực tiếp và không khó nếu $P$ là lớp các cấu xạ toàn ánh, tương
ứng radicial (các trường hợp này suy ra từ tính toàn ánh của $\alpha$), tương
ứng phẳng, tương ứng phẳng trung thành, tương ứng trơn (các trường hợp này suy
ra từ tính phẳng trung thành của $\alpha$), tương ứng kiểu hữu hạn. Dùng các
Định lý 1 và 2 cùng Bổ đề 1.2, ta cũng thấy điều tương tự đúng nếu $P$ là một
trong các lớp sau: các đẳng cấu, các phép nhúng mở, các phép nhúng đóng, các
phép nhúng (nếu $f$ kiểu hữu hạn và $Y$ Noether địa phương), các cấu xạ affine,
các cấu xạ hữu hạn, các cấu xạ quasi-hữu hạn, các cấu xạ mở, các cấu xạ đóng,
các đồng phôi, các cấu xạ tách,

\src{319}
các cấu xạ proper. Trường hợp quan trọng duy nhất chưa được giải quyết là các
cấu xạ xạ ảnh và quasi-xạ ảnh, đã được nêu trong nhận xét ở Mục 1.

\subsection*{3. Hạ xuống theo các cấu xạ hữu hạn phẳng trung thành}
\addcontentsline{toc}{subsection}{3. Hạ xuống theo các cấu xạ hữu hạn phẳng trung thành}

Cho $\alpha:S'\longrightarrow S$ là một cấu xạ hữu hạn tương ứng với một bó
đại số $\mathcal A'$ trên $S$, bó này, khi xem như bó môđun, là tự do địa phương
kiểu hữu hạn và là $\ne0$ ở mọi nơi; khi đó $\alpha$ là một cấu xạ phẳng trung
thành và quasi-compact, nên có thể áp dụng các kết quả trên. Việc cho một bó
tựa nhất quán $F'$ trên $S'$ tương đương với việc cho bó tựa nhất quán
$\alpha_*(F')$ trên $S$, được trang bị cấu trúc $\mathcal A'$-môđun của nó
(lưu ý rằng $\mathcal A'=\alpha_*(\mathcal O_{S'})$). Để đơn giản, bó này trên
$S$ cũng sẽ được ký hiệu là $F'$. Hai nghịch ảnh $p_i^*(F')$ của $F'$ trên
$S'\times_SS'$ cũng tương ứng với các bó tựa nhất quán gồm các
$(\mathcal A'\otimes_{\mathcal O_S}\mathcal A')$-môđun
\[
  F'\otimes_{\mathcal O_S}\mathcal A'
  \quad\text{và}\quad
  \mathcal A'\otimes_{\mathcal O_S}F'.
\]
Việc cho một $(S'\times_SS')$-đồng cấu từ bó thứ nhất vào bó thứ hai tương
đương với việc cho một đồng cấu của các
$(\mathcal A'\otimes_{\mathcal O_S}\mathcal A')$-môđun và, vì $\mathcal A'$
tự do địa phương, điều này tương đương với việc cho một đồng cấu của các
$(\mathcal A'\otimes\mathcal A')$-môđun:
\[
  \mathcal U=
  \operatorname{Hom}_{\mathcal O_S}(\mathcal A',\mathcal A')
  =\mathcal A'\otimes\check{\mathcal A}'
  \longrightarrow
  \operatorname{Hom}_{\mathcal O_S}(F',F'),
\]
nghĩa là, với mỗi tiết diện $\xi$ của $\mathcal U$ trên một tập mở $V$, cho một
đồng cấu $\mathcal O_S$-môđun
$T_\xi:F'|_V\longrightarrow F'|_V$ thỏa mãn các điều kiện
\begin{equation}
  T_{f\xi}(x)=fT_\xi(x),
  \qquad T_{\xi f}(x)=T_\xi(fx),
  \tag{3.1}
\end{equation}
trong đó $f,x$ lần lượt là các tiết diện của $\mathcal A',F'$ trên một tập mở
của $S$ nằm trong $V$. Khi ấy các điều kiện (i) và (ii) của một dữ liệu hạ xuống
(A, Mục 1, (c)) lần lượt được viết thành:
\begin{equation}
  T_{1_{\mathcal U}}x=x,
  \qquad\text{nghĩa là,}\qquad
  T_{1_{\mathcal U}}=\operatorname{id}_{F'},
  \tag{3.2}
\end{equation}
\begin{equation}
  T_{\xi\eta}=T_\xi T_\eta.
  \tag{3.3}
\end{equation}
Nói cách khác, \emph{một dữ liệu hạ xuống trên $F'$ tương đương với một biểu
diễn của bó $\mathcal O_S$-đại số
$\mathcal U=\operatorname{Hom}_{\mathcal O_S}(\mathcal A',\mathcal A')$ vào
bó $\mathcal O_S$-đại số
$\operatorname{Hom}_{\mathcal O_S}(F',F')$, thỏa mãn hai điều kiện tuyến tính
(3.1).} Nếu có một phép ghép cặp các bó tựa nhất quán trên $S'$:
\[
  F'_1\times F'_2\longrightarrow F'_3
\]
(có thể được hiểu như một phép ghép cặp các bó $\mathcal A'$-môđun trên $S$),

\src{320}
và các dữ liệu dán trên $F'_i$, được xác định bởi các đồng cấu
\[
  T_i:\mathcal U\longrightarrow
  \operatorname{Hom}_{\mathcal O_S}(F'_i,F'_i)
  \qquad (i=1,2,3),
\]
thì các dữ liệu ấy tương thích với phép ghép cặp đã cho, theo nghĩa hiển nhiên,
khi và chỉ khi điều kiện sau được thỏa mãn.

Với mọi tiết diện $\xi$ của $\mathcal U$ trên một tập mở, ký hiệu
\[
  \Delta\xi=\sum_i\xi'_i\otimes_{\mathcal A'}\xi''_i
\]
là tiết diện của $\mathcal U\otimes_{\mathcal A'}\mathcal U$ (trong đó
$\mathcal U$ được xem là một $\mathcal A'$-môđun qua cấu trúc môđun trái của
nó) được xác định bởi công thức
\[
  \xi\mathbin{\cdot}(fg)=\sum_i\xi'_i(f)\xi''_i(g)
\]
(trong đó $f$ và $g$ là hai tiết diện của $\mathcal A'$ trên một tập mở nhỏ
hơn), ta có công thức
\begin{equation}
  T^{(3)}_\xi(x\mathbin{\cdot}y)
  =\sum_iT^{(1)}_{\xi'_i}x\mathbin{\cdot}T^{(2)}_{\xi''_i}y
  \tag{3.4}
\end{equation}
với mọi cặp tiết diện $x,y$ của $\mathcal A'$ trên một tập mở nhỏ hơn. (Có
thể diễn đạt tính chất này bằng cách nói
rằng các đồng cấu $T^{(i)}$ tương thích với ánh xạ đường chéo của $\mathcal U$
đối với phép ghép cặp đã cho.) Đặc biệt, các công thức (3.1) đến (3.4) cho phép
ta diễn giải, theo ngôn ngữ các biểu diễn của đại số có ánh xạ đường chéo, các
dữ liệu hạ xuống trên một bó đại số tựa nhất quán trên $S'$, và do đó cũng diễn
giải (khi hạn chế vào các đại số giao hoán) các dữ liệu hạ xuống trên một
$S'$-lược đồ affine.

Từ đây ta chuyển sang một cách diễn giải tương tự cho các dữ liệu hạ xuống trên
một $S'$-tiền lược đồ $X'$ tùy ý: việc cho một $X'$ như vậy tương đương với
việc cho một tiền lược đồ $X'$ trên $S$, được trang bị một đồng cấu
$\mathcal O_S$-đại số
\[
  \mathcal A'\longrightarrow\mathcal O_{X'},
\]
và một dữ liệu hạ xuống trên $X'$ tương đương với việc cho một đồng cấu các bó
(tương thích với cấu xạ $h:X'\longrightarrow S'$)
\[
  \mathcal U\longrightarrow
  \operatorname{Hom}_{h^{-1}(\mathcal O_S)}
  (\mathcal O_{X'},\mathcal O_{X'}),
\]
thỏa mãn các điều kiện tương tự các điều kiện (3.1) đến (3.4) ở trên.

\textsc{Ví dụ 1 (``Weil'').} -- Giả sử $S'/S$ là một phủ étale Galois có nhóm
Galois $\Gamma$ (xem A, các Mục 3 và 4.d). Khi đó, một dữ liệu hạ xuống trên một
bó tựa nhất quán $F'$ trên $S'$ (tương ứng, trên một $S'$-tiền lược đồ $X'$)
tương đương với việc cho một biểu diễn của $\Gamma$ bằng các tự đẳng cấu của
$(S',F')$ (tương ứng, của $(S',X')$), tương thích với tác động của $\Gamma$
trên $S'$. Kết quả này

\src{321}
là ``hình thức'', nghĩa là được chứng minh theo ngôn ngữ phạm trù; nhưng theo
quan điểm của mục này, nó cũng suy ra từ cấu trúc tường minh của $\mathcal U$
(được trang bị cấu trúc vành, đồng cấu vành
$\mathcal A'\longrightarrow\mathcal U$ và ánh xạ đường chéo), cấu trúc này được
biết hoàn toàn nhờ kết quả sau: $\mathcal U$, khi xem như một
$\mathcal A'$-môđun trái, có một cơ sở gồm các tiết diện của $\mathcal U$
tương ứng với các phần tử của $\Gamma$.

\textsc{Ví dụ 2 (``Cartier'').} -- Cho $p$ là một số nguyên tố; giả sử
$p\mathcal O_S=0$ (nghĩa là $\mathcal O_S$ có đặc số $p$),
$(\mathcal A')^p\subset\mathcal O_S=\mathcal A$ (nghĩa là $S'/S$ radicial có
chiều cao 1), và bó đại số $\mathcal A'$ trên $\mathcal A$ có một $p$-cơ sở
tại địa phương, nghĩa là một họ tiết diện $(x_i)$ sao cho $\mathcal A'$ được sinh
như một đại số bởi các $x_i$, chỉ chịu các điều kiện $x_i^p=0$. Ta giả sử tập
các chỉ số $i$ là hữu hạn, có lực lượng $n$. Cho $\mathfrak g$ là bó các
$\mathcal A$-đạo hàm của $\mathcal A'$; đó là một bó $\mathcal A'$-môđun tự do
địa phương hạng $n$, và hơn nữa là một bó $p$-đại số Lie trên $\mathcal A$
(nhưng không trên $\mathcal A'$), thỏa mãn điều kiện
\begin{equation}
  [X,fY]=X(f)Y+f[X,Y].
  \tag{3.5}
\end{equation}

\textsc{Bổ đề.} -- $\mathcal U=
\operatorname{Hom}_{\mathcal O_S}(\mathcal A',\mathcal A')$ được sinh, như một
$\mathcal O_S$-đại số được trang bị một đồng cấu đại số
$\mathcal A'\longrightarrow\mathcal U$, bởi $\mathcal A'$-môđun con trái
$\mathfrak g$, với các quan hệ bổ sung
\begin{equation}
  \left\{
  \begin{aligned}
    Xf-fX&=X(f),\\
    XY-YX&=[X,Y],\\
    X^p&=X^{(p)}.
  \end{aligned}
  \right.
  \tag{3.6}
\end{equation}

Từ bổ đề trên suy ra rằng một dữ liệu hạ xuống trên bó tựa nhất quán $F'$ trên
$S'$ tương đương với việc cho, với mọi $X\in\mathfrak g$, một tự đồng cấu
$\mathcal O_S$-môđun $\overline X$ của $F'$ thỏa mãn các điều kiện
\begin{align}
  \overline{fX}&=f\overline X, \tag{3.7}\\
  \overline X(fx)&=X(f)x+f\overline X(x), \tag{3.8}\\
  \overline{[X,Y]}&=[\overline X,\overline Y], \tag{3.9}\\
  \overline{X^{(p)}}&=\overline X^p. \tag{3.10}
\end{align}
(Đó là điều có thể gọi là một liên thông tuyến tính trên $F'$, không có độ
cong và tương thích với lũy thừa bậc $p$.) Khái niệm dữ liệu hạ xuống trên một
$S'$-tiền lược đồ $X'$ cũng được làm tường minh theo cách ấy; ở đây quan hệ
(3.4) được thay bằng điều kiện rằng các $\overline X$ là các đạo hàm của
$\mathcal O_{X'}$. Vì cấu xạ $S'\longrightarrow S$ là radicial, Định lý 3 bảo
đảm rằng mọi dữ liệu hạ xuống như vậy đều

\src{322}
hữu hiệu và do đó xác định một $S$-tiền lược đồ $X$.

Ta lưu ý rằng không cần đặt bất kỳ giả thiết nào về tính phẳng, tính không kỳ
dị hay tính hữu hạn đối với $F'$, tương ứng $X'$.

\subsection*{4. Ứng dụng vào các tiêu chuẩn về tính hữu tỉ}
\addcontentsline{toc}{subsection}{4. Ứng dụng vào các tiêu chuẩn về tính hữu tỉ}

Cho $X$ là một $S$-tiền lược đồ sao cho ảnh trực tiếp của $\mathcal O_X$ trên
$S$ là $\mathcal O_S$; tính chất này vẫn đúng sau mọi mở rộng phẳng
$S'\longrightarrow S$ của cơ sở $S$. Nếu $F$ là một bó khả nghịch (nghĩa là
tự do địa phương hạng 1) trên $X$, thì các tự đẳng cấu của $F$, được đồng nhất
với các tiết diện khả nghịch của $\mathcal O_X$, tương ứng song ánh với các
tiết diện khả nghịch của $\mathcal O_S$. Bây giờ cho $s$ là một tiết diện của
$X$ trên $S$; theo cách nói gợi hình, ta sẽ gọi là một tiết diện của $F$ trên
$s$ một tiết diện của bó khả nghịch $s^*(F)$ trên $S$. Từ những điều trên suy ra
rằng, nếu $F_i$ ($i=1,2$) là hai bó khả nghịch trên $X$, mỗi bó được trang bị
một tiết diện trên $s$, và nếu $F_1$ và $F_2$ đẳng cấu, thì tồn tại duy nhất
một đẳng cấu từ $F_1$ lên $F_2$ tương thích với các tiết diện đang xét (nghĩa
là biến tiết diện thứ nhất thành tiết diện thứ hai). Mặt khác, không phụ thuộc
vào tiết diện $s$, ta quy ước coi hai bó khả nghịch $F_1$ và $F_2$ trên $X$ là
tương đương nếu mọi điểm của $S$ đều có một lân cận mở $U$ sao cho các hạn chế
của $F_1$ và $F_2$ lên $X|U$ đẳng cấu. Khi đó mọi bó khả nghịch $F$ trên $X$
tương đương với một bó khả nghịch $F_1$ được trang bị một tiết diện được đánh
dấu trên $s$ (lấy $F_1=F\,s^*(F)^{-1}$), và $F_1$ được xác định đến đẳng cấu.
Nói cách khác, việc phân loại các bó khả nghịch trên $X$ đến tương đương cũng
chính là việc phân loại, đến đẳng cấu, các bó khả nghịch được trang bị một tiết
diện được đánh dấu.

Vì các tính chất này vẫn đúng sau một mở rộng phẳng
$\alpha:S'\longrightarrow S$ của cơ sở (thay tiết diện $s$ bằng ảnh ngược $s'$
của nó qua $\alpha$), nên từ Định lý 1 ta suy ra:

\emph{Cho tiền lược đồ $X/S$ như trên và giả sử nó có một tiết diện $s$; cho
$\alpha:S'\longrightarrow S$ là một cấu xạ phẳng trung thành và quasi-compact;
cho $F'$ là một bó khả nghịch trên $X'=X\times_SS'$. Điều kiện cần và đủ để
$F'$ tương đương với ảnh ngược trên $X'$ của một bó khả nghịch $F$ trên $X$ là
hai ảnh ngược $p_1^*(F')$ và $p_2^*(F')$ của nó trên
$X'\times_XX'=X\times_S(S'\times_SS')$ tương đương. Khi đó $F$ được xác định
đến tương đương.} (Khi ấy ta nói $F'$ hữu tỉ trên $S$.)

Áp dụng nguyên lý này khi $\alpha:S'\longrightarrow S$ như trong Ví dụ 1 hoặc
Ví dụ 2 của mục trước, ta thu được các tiêu chuẩn hữu tỉ của Weil hoặc Cartier.
(Lưu ý rằng các tác giả này chỉ xét trường hợp $S$ và $S'$

\src{323}
là phổ của các trường; do đó, hiển nhiên $S$ là phổ của một vành địa phương,
và quan hệ tương đương được đưa vào ở trên chính là quan hệ đẳng cấu.) Trong
trường hợp thứ nhất, $F'$ hữu tỉ trên $S$ khi và chỉ khi các biến đổi của nó
dưới $\Gamma$ tương đương với $F'$. Để phát biểu tiêu chuẩn hữu tỉ trong trường
hợp thứ hai, nói chung ta xét cấu xạ đường chéo
\[
  X'\longrightarrow X''=X'\times_XX'
\]
của $X'/X$, bó iđêan tương ứng $\mathcal I$ trên $X'\times_XX'$, và bó
$\mathcal I/\mathcal I^2$, được đồng nhất với ảnh ngược
$\Omega^1_{X'/X}$ của nó trên $X'$. Vì các hạn chế lên đường chéo của
$F''_i=p_i^*(F')$ ($i=1,2$) đẳng cấu (do cả hai đều đẳng cấu với $F'$), nghĩa
là $F''_1(F''_2)^{-1}=F''$ có hạn chế tầm thường trên đường chéo, nên hạn chế
của $F''$ lên $(X'',\mathcal O_{X''}/\mathcal I^2)$ được cho, đến đẳng cấu,
bởi một phần tử xác định duy nhất $\xi$ của
\[
  H^1(X'',\mathcal I/\mathcal I^2)
  =H^1(X',\Omega^1_{X'/X}).
\]
Hơn nữa, trong trường hợp này ta có
\[
  \Omega^1_{X'/X}
  =\Omega^1_{S'/S}\otimes_{\mathcal O_S}\mathcal O_X,
\]
và do đó, nếu $\Omega^1_{S'/S}$ tự do địa phương trên $S$ (như trong trường
hợp Cartier), thì $\xi$ xác định một tiết diện của
\[
  R^1f'_*(\mathcal O_{X'})\otimes\Omega^1_{S'/S}
\]
trên $S'$ (gọi là lớp Atiyah--Cartier của bó khả nghịch $F'$ trên $X'/S$).
Sự triệt tiêu của lớp này là điều kiện cần và đủ để các ảnh ngược của $F'$ qua
hai phép chiếu từ
\[
  (X'',\mathcal O_{X''}/\mathcal I^2)
  =X\times_S(S'',\mathcal O_{S''}/\mathcal J^2)
\]
lên $X'$ tương đương (trong đó $\mathcal J$ là bó iđêan trên
$S''=S'\times_SS'$ được xác định bởi cấu xạ đường chéo
$S'\longrightarrow S'\times_SS'$). Vì vậy, sự triệt tiêu này hiển nhiên là
điều kiện cần để các ảnh ngược của $F'$ trên chính
$X''=X\times_SS''$ tương đương, và do đó cũng là điều kiện cần để $F'$ tương
đương với ảnh ngược của một bó khả nghịch $F$ trên $X$. Lớp Atiyah--Cartier
cũng có thể được hiểu là cản trở đối với sự tồn tại, cục bộ trên $S'$, của một
liên thông trên $F'$ tương đối với các đạo hàm của $X'/X$; hơn nữa, một liên
thông như vậy được xác định khi đã biết các đạo hàm của $F'$ tương ứng với các
mở rộng tự nhiên lên $X'$ của các đạo hàm của $S'/S$. Từ điều này và các khai
triển trong mục trước, ta dễ dàng suy ra rằng trong trường hợp Ví dụ 2 của mục
ấy, khi $X/S$ có một tiết diện, sự triệt tiêu của lớp Atiyah--Cartier cũng đủ
để $F'$ hữu tỉ trên $S$.

\subsection*{5. Ứng dụng vào phép thu hẹp cơ sở trong một lược đồ abel}
\addcontentsline{toc}{subsection}{5. Thu hẹp cơ sở trong một lược đồ abel}

Cho $S$ là một tiền lược đồ. Một lược đồ abel trên $S$ là một lược đồ $X$ trơn
và proper

\src{324}
trên $S$ sao cho các thớ tại những điểm $x\in S$ là các lược đồ của những đa
tạp abel trên $\kappa(x)$. Giả sử $S$ Noether và chính quy (nghĩa là các vành
địa phương của nó chính quy). Khi đó, dùng định lý liên thông của Murre [4]
(ít nhất trong trường hợp ``đồng đặc số'', là trường hợp mà định lý được dẫn
hiện đã được chứng minh), ta có thể chứng minh rằng mọi tiết diện hữu tỉ của
$X$ trên $S$ đều xác định khắp nơi (nghĩa là thực sự là một tiết diện); điều
này tổng quát hóa một định lý cổ điển của Weil. Tổng quát hơn, từ đó suy ra
rằng nếu $X'$ là một lược đồ trơn trên $S$, thì mọi $S$-cấu xạ hữu tỉ từ $X'$
vào $X$ đều xác định khắp nơi. Ta thu được kết quả sau, tổng quát hóa một kết
quả của Igusa--Lang: Cho $S$ Noether và chính quy, và ký hiệu $K$ là
vành các hàm hữu tỉ của nó (một tích trực tiếp các trường); cho $X$ là một lược
đồ abel trên $K$. Nếu $X$ đẳng cấu với một $K$-lược đồ có dạng
$X_0\times_S\operatorname{Spec}(K)$, trong đó $X_0$ là một lược đồ abel trên
$S$, thì $X_0$ được xác định đến một đẳng cấu duy nhất.

Dùng kết quả duy nhất ở trên, ta thấy rằng bài toán thu hẹp cơ sở trong $X$ là
cục bộ trên $S$ (và do đó chỉ cần biết cách thực hiện phép thu hẹp trên các
$\operatorname{Spec}(\mathcal O_x)$, với $x\in S$). Tương tự, nếu
$\alpha:S'\longrightarrow S$ là một cấu xạ trơn kiểu hữu hạn, nếu $K'$ là
vành các hàm hữu tỉ của $S'$, và nếu $X\otimes_KK'$ có dạng
$X'_0\times_{S'}\operatorname{Spec}(K')$, thì $X'_0$ được trang bị một dữ liệu
hạ xuống chính tắc tương đối với $\alpha$. Theo Định lý 3, ta suy ra:

\result{MỆNH ĐỀ}{5.1} Cho $S$ là một tiền lược đồ Noether, chính quy và bất khả
quy, có trường hàm hữu tỉ $K$; cho $K'$ là một mở rộng hữu hạn của $K$, không
ram hóa trên $S$; cho $S'$ là chuẩn hóa của $S$ trong $K'$ (do đó là một phủ
étale của $S$); và cho $X$ là một lược đồ abel trên $K$ sao cho
$X\otimes_KK'$ có dạng $X'_0\times_{S'}\operatorname{Spec}(K')$, trong đó
$X'_0$ là một lược đồ abel xạ ảnh trên $S'$. Khi đó $X$ có dạng
$X_0\times_S\operatorname{Spec}(K)$, trong đó $X_0$ là một lược đồ abel xạ ảnh
trên $S$.

\result{NHẬN XÉT}{} Người trình bày không biết liệu có thể thay giả thiết
$S'\longrightarrow S$ là một phủ étale toàn ánh (cho phép dùng Định lý 3) bằng
giả thiết rằng đó là một cấu xạ trơn, toàn ánh và kiểu hữu hạn (ngay cả khi giả
sử nó là étale), hoặc liệu mệnh đề có còn đúng mà không cần giả sử $X'_0$ xạ
ảnh trên $S'$ hay không (một điều kiện có lẽ tự động được thỏa mãn).

\subsection*{6. Ứng dụng vào các tiêu chuẩn tầm thường cục bộ và đẳng tầm}
\addcontentsline{toc}{subsection}{6. Các tiêu chuẩn tầm thường cục bộ và đẳng tầm}

Cho $S$ là một tiền lược đồ, $G$ là một ``tiền lược đồ nhóm'' trên $S$, và $P$
là một tiền lược đồ trên $S$ mà trên đó ``$G$ tác động'' (bên phải). Ta nói
$P$ là thuần nhất chính hình thức dưới $G$ nếu cấu xạ quen thuộc
\[
  G\times_SP\longrightarrow P\times_SP
\]

\src{325}
cảm sinh bởi tác động của $G$ lên $P$ là một đẳng cấu. Từ đây trở đi, ta giả
sử $G$ phẳng trên $S$ (và do đó phẳng trung thành trên $S$), và dành tên
phân thớ chính thuần nhất dưới $G$ cho một $P$ thuần nhất chính hình thức
mà phẳng trung thành và quasi-compact trên $S$. Điều này hiển nhiên tương đương
với việc tồn tại một mở rộng phẳng trung thành và quasi-compact
$S'\longrightarrow S$ của cơ sở $S$ sao cho phân thớ thuần nhất chính hình thức
\[
  P'=P\times_SS'
\]
dưới $G'=G\times_SS'$ là tầm thường, nghĩa là đẳng cấu với $G'$ (nghĩa là có
một tiết diện); đặc biệt, có thể lấy $S'=P$. Ngoài ra, lưu ý rằng nếu $S$ địa
phương Noether, thì giả thiết $P$ phẳng trung thành tương đương với
giả thiết
\[
  \widehat P_s=P\times_S\operatorname{Spec}(\widehat{\mathcal O}_s)
\]
phẳng trung thành trên $\widehat{\mathcal O}_s$ với mọi $s\in S$ (trong đó
$\widehat{\mathcal O}_s$ ký hiệu phép đầy đủ hóa của vành địa phương
$\mathcal O_s$), vì $\widehat{\mathcal O}_s$ phẳng trung thành trên
$\mathcal O_s$. Hơn nữa, nếu $P$ có kiểu hữu hạn trên $S$ địa phương Noether,
thì tập các điểm $s$ thỏa điều kiện trên là dựng được. Vì vậy, nếu $S$ là một
``tiền lược đồ Jacobson'' (chẳng hạn một lược đồ kiểu hữu hạn trên một trường,
hoặc tổng quát hơn là một vành Jacobson), thì chỉ cần kiểm tra điều kiện
đang xét tại các điểm đóng của $S$. Điều này quy ta về trường hợp cơ sở là phổ
của một vành địa phương đầy đủ $A$. Khi $S=\operatorname{Spec}(A)$ (với $A$ là
một vành địa phương Noether đầy đủ) và $P$ có kiểu hữu hạn trên $S$, tính phẳng
trung thành của $P/S$ cũng tương đương với sự tồn tại của một $S'$ hữu hạn và
phẳng trên $S$ sao cho $P'$ tầm thường; nếu thêm $G$ trơn trên $S$, thì có thể
giả sử $S'$ étale trên $S$. Do đó, nếu trường thặng dư của $A$ còn đóng đại số
(``trường hợp hình học''), thì $P$ phẳng trung thành trên $A$ khi và chỉ khi nó
tầm thường. Như vậy, nếu $S$ là một tiền lược đồ đại số trên một trường đóng
đại số và $G$ trơn, kiểu hữu hạn trên $S$, thì điều kiện phẳng trung thành trên
$S$ tương đương với điều kiện tầm thường giải tích (SLF) của Serre ([6],
tr.~1--12).

Ta có thể đưa vào những loại điều kiện khác, mạnh hơn, trên $P$ và mang bản chất
của một ``tính tầm thường cục bộ''. Cụ thể, ta nói $P$ là đẳng tầm
(tương ứng, đẳng tầm nghiêm ngặt) nếu với mọi $s\in S$, tồn tại một lân
cận mở $U$ của $s$ và một cấu xạ hữu hạn, phẳng trung thành (tương ứng, một phủ
étale toàn ánh) $U'\longrightarrow U$ sao cho $P'=P\times_SU'$ tầm thường.
(Thuật ngữ của chúng ta khác với Serre [6], người gọi đẳng tầm cục bộ
điều mà chúng ta gọi là đẳng tầm nghiêm ngặt.) Tính đẳng tầm nghiêm ngặt
đặc biệt hữu ích khi $G$ trơn trên $S$, nhưng ngược lại là một khái niệm không
thích hợp trong những trường hợp khác.

Nếu $G$ affine trên $S$, thì theo đoạn 2, mọi phân thớ chính thuần nhất $P$ dưới
$G$ đều affine; bởi vậy, nhờ Định lý 2, ta có thể ``hạ xuống''

\src{326}
các phân thớ như vậy dọc theo những cấu xạ phẳng trung thành và quasi-compact.
Đặc biệt, lấy $G=\operatorname{Gl}(n)_S$, được xác định bởi điều kiện rằng hàm
tử trên các $S$-tiền lược đồ
\[
  S'\longmapsto\operatorname{Hom}_S(S',G)
\]
(nhận giá trị trong phạm trù các nhóm) được đồng nhất với hàm tử
\[
  \operatorname{Gl}(n)(S')
  =\operatorname{Gl}\bigl(n,H^0(S',\mathcal O_{S'})\bigr)
\]
trong A, đoạn 4.e. Dùng các sự kiện (i) mọi phân thớ chính thuần nhất dưới $G$
(tương ứng, mọi bó tự do địa phương hạng $n$ trên $S$) trở thành đẳng cấu với
đối tượng ``tầm thường'' $G$ (tương ứng, $\mathcal O_S^n$) sau một mở rộng
phẳng trung thành và quasi-compact thích hợp của $S$; (ii) các đối tượng thuộc
loại đang xét (các phân thớ chính thuần nhất dưới $G$, tương ứng các bó tự do
địa phương hạng $n$) có thể được hạ xuống dọc theo những cấu xạ như vậy; và
cuối cùng, (iii) nhóm tự đẳng cấu của phân thớ tầm thường trên một $S'/S$ đẳng
cấu một cách hàm tử với nhóm tự đẳng cấu của bó tự do địa phương tầm thường
hạng $n$ trên $S'$, ta suy ra một cách ``hình thức'' rằng việc cho trên $S$
(hoặc trên một $S'/S$) một phân thớ chính thuần nhất có nhóm $G$ ``cũng chính
là'' việc cho tại đó một bó tự do địa phương hạng $n$. (Chính xác hơn, ta có
một tương đương giữa các phạm trù phân thớ.) Đặc biệt, ta suy ra:

\result{MỆNH ĐỀ}{6.1} Mọi phân thớ chính thuần nhất có nhóm
$\operatorname{Gl}(n)_S$ đều tầm thường cục bộ.

Bằng những lập luận quen thuộc, ta thu được cùng kết quả cho các nhóm cấu trúc
như $\operatorname{Sl}(n)_S$, $\operatorname{Sp}(n)_S$ và các tích của những
nhóm như vậy. Ta cũng suy ra rằng, nếu $F$ là một nhóm con đóng của
$G=\operatorname{Gl}(n)_S$, phẳng trên $S$, sao cho thương $G/F$ tồn tại, và
nếu $G$ là một phân thớ chính thuần nhất đẳng tầm (tương ứng, đẳng tầm nghiêm
ngặt) trên $G/F$, với nhóm $F\times_S(G/F)$, thì mọi phân thớ chính thuần nhất
có nhóm $F$ đều đẳng tầm (tương ứng, đẳng tầm nghiêm ngặt). Điều này áp dụng
cho tất cả các ``nhóm tuyến tính'' trên $S$ đã được dùng cho đến đây, và đặc
biệt khi $G=S\times_k\Gamma$, trong đó $S$ là một tiền lược đồ trên trường $k$
và $\Gamma$ là một nhóm tuyến tính theo nghĩa cổ điển, đặc biệt là một nhóm
trơn, trên $k$. Do đó, đối với các nhóm như vậy, điều này giải quyết một câu hỏi
của Serre (loc. cit.).

Ta cũng lưu ý rằng, đối với phần lớn các nhóm trơn trên $S$ đã biết, dù tuyến
tính hay không, và trong mọi trường hợp đối với các nhóm có dạng
$S\times_k\Gamma$ như trên, có thể chứng minh rằng mọi phân thớ chính thuần nhất
đẳng tầm đều đẳng tầm nghiêm ngặt. Điều này đặc biệt giải quyết một câu hỏi khác
của Serre (loc. cit., tr.~1--14), vì một phân thớ chính thuần nhất thu được bằng
phép hạ xuống kiểu Cartier (xem đoạn 3, Ví dụ 2) hiển nhiên là đẳng tầm.

\result{NHẬN XÉT}{} Một trong những khó khăn cốt yếu của các vấn đề này (ngoài
vấn đề về sự tồn tại của các lược đồ thương $G/F$) là thiếu những tiêu chuẩn về
tính hữu hiệu cho một dữ liệu hạ xuống dọc theo một cấu xạ phẳng trung thành
không hữu hạn.

\src{327}
\begin{center}
  \textbf{TÀI LIỆU THAM KHẢO}
\end{center}

\begin{description}
  \item[{[1]}] Dieudonné (J.) và Grothendieck (Alexander). --- \emph{Éléments
    de géométrie algébrique}, sắp xuất bản trong Publications mathématiques de
    l'Institut des Hautes Études Scientifiques.
  \item[{[2]}] Grauert (H.) und Remmert (R.). --- \emph{Komplexe Räume}, Math.
    Annalen, tập 136, 1958, tr.~245--318.
  \item[{[3]}] Grothendieck (Alexander). --- \emph{Géométrie formelle et
    géométrie algébrique}. Séminaire Bourbaki, tập 11, 1958/59, số 182.
  \item[{[4]}] Murre (J. P.). --- \emph{On a connectedness theorem for a
    birational transformation at a simple point}, Amer. J. Math., tập 80, 1958,
    tr.~3--15.
  \item[{[5]}] Serre (Jean-Pierre). --- \emph{Géométrie algébrique et géométrie
    analytique}, Ann. Inst. Fourier Grenoble, tập 6, 1955--56, tr.~1--42.
  \item[{[6]}] Serre (Jean-Pierre). --- \emph{Espaces fibrés algébriques},
    Séminaire Chevalley, tập 2, 1958: Anneaux de Chow et applications, số 1.
\end{description}

\par
